\documentclass[11pt]{article}
\usepackage[T1]{fontenc}
\usepackage{lmodern}
\usepackage[margin=1in]{geometry}
\usepackage{amsmath,amssymb,amsthm,mathtools}
\usepackage{booktabs,array,longtable}
\usepackage{enumitem}
\usepackage{xcolor}
\usepackage{microtype}
\usepackage[colorlinks=true,linkcolor=blue!55!black,citecolor=blue!55!black,urlcolor=blue!55!black]{hyperref}
\usepackage[nameinlink,noabbrev]{cleveref}

\hypersetup{
  pdftitle={Rank-Three Mapping-Class-Finite Representations in Genus Four: A Candidate Extension},
  pdfauthor={Robert Sneiderman},
  pdfsubject={New, unrefereed candidate extension},
  pdfkeywords={mapping class groups, surface groups, local systems, genus four, representation theory}
}

\newtheorem{theorem}{Theorem}[section]
\newtheorem{proposition}[theorem]{Proposition}
\newtheorem{lemma}[theorem]{Lemma}
\newtheorem{corollary}[theorem]{Corollary}
\newtheorem{candidate}[theorem]{Candidate theorem}
\theoremstyle{definition}
\newtheorem{definition}[theorem]{Definition}
\newtheorem{remark}[theorem]{Remark}
\newtheorem{audit}[theorem]{Audit point}

\newcommand{\C}{\mathbf C}
\newcommand{\Ocal}{\mathcal O}
\newcommand{\Mcal}{\mathcal M}
\newcommand{\Vcal}{\mathcal V}
\newcommand{\Ucal}{\mathcal U}
\newcommand{\Mod}{\operatorname{Mod}}
\newcommand{\rk}{\operatorname{rk}}
\newcommand{\degp}{\operatorname{pardeg}}
\newcommand{\ad}{\operatorname{ad}}
\newcommand{\im}{\operatorname{im}}
\newcommand{\End}{\operatorname{End}}
\newcommand{\tr}{\operatorname{tr}}
\newcommand{\PGL}{\operatorname{PGL}}
\newcommand{\GL}{\operatorname{GL}}
\newcommand{\SL}{\operatorname{SL}}
\newcommand{\SO}{\operatorname{SO}}
\newcommand{\Sp}{\operatorname{Sp}}
\providecommand{\tightlist}{\setlength{\itemsep}{0pt}\setlength{\parskip}{0pt}}

\title{Rank-Three Mapping-Class-Finite Representations\\
in Genus Four: A Candidate Extension}
\author{Robert Sneiderman}
\date{Version 0.1.1-candidate\\15 August 2026}

\begin{document}
\maketitle

\begin{abstract}
We give a candidate extension to genus four of the rank-three
finite-image argument developed in the companion genus-at-least-five
candidate manuscript.  For an irreducible rank-three system, the new work
eliminates the dense rank-eight adjoint coefficient, the rank-five
orthogonal coefficient, and the rank-six monomial coefficient on every
finite-etale moduli cover.  It then resolves the scalar equality case in
formal propagation by comparing the adjoint Killing form with the standard
rank-eight symplectic monodromy representation.  The conclusion remains
conditional on the fixed-part, Hodge, parabolic, canonical-Deligne,
finite-cover, integrality, isomonodromy, and socle interfaces stated below.
The argument is new and unrefereed, and is presented as a
\texttt{CANDIDATE}, not as an established theorem.
\end{abstract}

\begin{center}
\fbox{\parbox{0.92\textwidth}{
\textbf{Status and scope.}  Every result called a ``candidate theorem'' in
this manuscript is conditional on the imported interface package in
\cref{sec:interfaces}.  The companion genus-at-least-five manuscript is
itself an unrefereed candidate, not a published input.  This paper is
restricted to genus four and makes no lower-genus assertion.}}
\end{center}

\begin{center}
\small Copyright \textcopyright\ 2026 Robert Sneiderman.  Licensed under
\href{https://creativecommons.org/licenses/by/4.0/}{CC BY 4.0}.
\end{center}

\tableofcontents

\section{Statement and proof architecture}\label{sec:statement}

Let $\Sigma_{4,n}$ be an orientable genus-four surface with $n$ punctures.
The mapping class group $\Mod_{4,n}$ acts on conjugacy classes of
representations of $\pi_1(\Sigma_{4,n})$ by precomposition.

\begin{candidate}[Rank-three finite image in genus four]\label{cand:main}
Assume the published inputs and the companion-candidate interfaces recorded
in \cref{sec:interfaces}.  Let $n\ge0$.  If
\[
 \rho:\pi_1(\Sigma_{4,n})\longrightarrow\GL_3(\C)
\]
has finite $\Mod_{4,n}$-orbit up to conjugacy, then $\rho$ has finite image.
\end{candidate}

Landesman--Litt prove finite image under the strict inequality
$r<\sqrt{g+1}$ \cite[Theorem~1.2.1]{LLCan}.  Their published theorem therefore
does not directly treat $(g,r)=(4,3)$.  The companion manuscript proposes a
rank-three theorem for $g\ge5$ \cite{SneidermanG5}; the present paper is a
separate extension, and does not promote either candidate to established
status.

The proof has five layers.
\begin{enumerate}[label=\textup{(\arabic*)}]
\item The fixed-part reduction leaves three coefficient ranks: the dense
rank-eight adjoint coefficient, a rank-five orthogonal coefficient, and a
rank-six monomial coefficient.
\item Exact no-high and high-Harder--Narasimhan ledgers are reconstructed in
genus four.  Universal-extension symmetry, normalizer cocycles, Killing and
Spin parity, canonical-lattice pairing bounds, and low-degree section
estimates eliminate every dense rank-eight branch.
\item The rank-five coefficient is eliminated by the first-super endpoint
classification and its invariant symmetric form.
\item Exact monomial arithmetic leaves two no-high rows and one strict
high-HN row for the rank-six coefficient.  Symmetric universal extensions,
parabolic stability, and the determinant-sign character eliminate them.
\item Projective finite-index descent and Artin devissage are unchanged.
The only new propagation equality identifies a putative rank-eight adjoint
coefficient with the standard genus-four cohomology system; their preserved
bilinear forms are incompatible.
\end{enumerate}

\section{Published inputs and imported candidate interfaces}
\label{sec:interfaces}

The dependency boundary is part of the statement.  We use the following
published results in the forms already isolated in the companion manuscript:
\begin{itemize}
\item Landesman--Litt's projective globalization, fixed-part, Hodge,
parabolic multiplication, rigidity, integrality, isomonodromy, and
characteristic-socle interfaces \cite{LLCan,LLGeo};
\item the small-slope vector-bundle Clifford and generated-subbundle
estimates of Brambila-Paz--Grzegorczyk--Newstead \cite{BPGN};
\item the classification of holomorphic sections of universal curves by
Earle--Kra \cite{EarleKra}, and Chen--Salter's virtual non-splitting theorem
in its stated range $g\ge4$ \cite{ChenSalter}.
\end{itemize}

The following items are imported from the unrefereed companion candidate
\cite{SneidermanG5}.  They are hypotheses of Candidate
theorem~\ref{cand:main}, not claims of published provenance:
\begin{enumerate}[label=\textup{B\arabic*},leftmargin=3.2em]
\item the passage from an MCG-finite rank-three system to its actual
multiplicity-free adjoint constituents on connected finite-etale moduli
covers;
\item the canonical-Deligne and coparabolic convention, the fixed-part
Hodge-line reduction, and the closed-descent interface;
\item the strict endpoint and first-super endpoint estimates, including the
weight-one universal-extension classification;
\item the normalizer boundary-cocycle obstruction and the extension
compatibility identities for bracket and invariant symmetric forms;
\item projective finite-index descent, direct Artin devissage, the exact
trivial-coefficient gap, reducible unitary multiplicity control, and the
remaining projective/linear integrality, unitary-conjugate, and
characteristic-socle propagation steps.
\end{enumerate}

For clarity, the imported projective-closure reduction is
\[
\begin{array}{c|c}
\text{projective closure}&\ad V\\ \hline
\PGL_3\text{-dense}&U_8\\
\PGL_2\simeq\SO_3&U_5\oplus U_3\\
\text{monomial, component quotient }S_3&U_6\oplus U_2\\
\text{monomial, component quotient }C_3&
 U_3^+\oplus U_3^-\oplus\chi\oplus\chi^{-1}.
\end{array}
\]
Ranks at most three are already below the relevant fixed-part threshold in
genus four.  The new fixed-part work is therefore exactly the elimination of
$U_8$, $U_6$, and $U_5$ carried out below.

\begin{audit}[Conditionality]
Finite HN enumeration and fibrewise linear algebra do not certify the
geometric interfaces in B1--B5.  In particular, compilation of this paper or
successful execution of finite verification scripts is not a proof of
Candidate theorem~\ref{cand:main}.
\end{audit}

\section{Genus-four setup}\label{sec:setup}

Work on a very general smooth genus-four fibre $C$ of the relevant finite
cover.  It is nonhyperelliptic.  For a parabolically stable coefficient
$E_*$ of parabolic degree zero, let $E$ denote the underlying canonical
Deligne bundle and let $S=-\deg E$ be its total canonical parabolic weight.
For a rank-zero Hodge vector, let $M\subset E$ be the saturated Hodge line,
put $L=M^{-1}$ and $G=E/M$, and let $A\subset E^\vee\otimes\omega_C$ be the
saturated evaluation kernel.  In genus four set
\[
 q=\deg(M^\vee\otimes\omega_C)=6-\deg M.
\]
When the high-HN part $N\subset A$ is nonzero, write
$t=\rk N$ and $y=\deg N-S$.  All saturations and brackets below are taken in
the canonical parahoric lattice supplied by B2 and B4.

The next two sections give the exhaustive numerical ledgers for the dense
rank-eight coefficient.  Subsequent sections eliminate every listed row.

\section{No-high reconstruction}\label{no-high-reconstruction}

For the rank-eight dense adjoint coefficient, the no-high inequalities
give

\[S\le \frac{16(8-g)}9,
\qquad
\max(2g-1,2g-2+S)\le q\le14-S.\]

The quotient bundle has

\[(\mathrm{rk}\,G,\deg G,h^0(G))
=(7,q-(2g-2)-S,q-g+1).\]

The exhaustive HN enumeration imposes strictly decreasing HN slopes. For
\(S=0\) it uses positive-degree factors, as in the companion manuscript;
for \(S>0\) it also
permits degree-zero factors, and every factor has slope at most
\(2g-2\). The section bounds are applied on the very general,
nonhyperelliptic genus-four fibre.

\subsection{Numerical survivor table}\label{genus-four}

\begin{longtable}[]{@{}rlll@{}}
\toprule\noalign{}
\(S\) & Candidate \(q\) & Eliminated numerically & Survivors \\
\midrule\noalign{}
\endhead
\bottomrule\noalign{}
\endlastfoot
0 & 7--14 & 12, 13 & 7, 8, 9, 10, 11, 14 \\
2 & 8--12 & none & 8, 9, 10, 11, 12 \\
3 & 9--11 & 11 & 9, 10 \\
4 & 10 & none & 10 \\
\end{longtable}

Thus 14 of 17 candidates survive. The first is

\[(S,q)=(0,7),
\qquad
(\mathrm{rk}\,G,\deg G,h^0(G))=(7,1,4).\]

The exact BPGN small-slope inequality gives only \(h^0(G)\le5\) for a
semistable rank-seven, degree-one piece in genus four. Requiring four
sections is therefore genuinely compatible with the cited bound. Section
counting cannot remove this branch.

\section{High-HN numerical frontier}\label{high-hn-numerical-frontier}

Write \(t=\mathrm{rk}\,N\) and \(y=\deg N-S\). After the offset,
Clifford, lower-slope, and \(q\) inequalities, the following necessary
rows remain. They have not been geometrically realized as adjoint
branches.

\subsection{Numerical survivor table}\label{genus-four-1}

\begin{longtable}[]{@{}rrll@{}}
\toprule\noalign{}
\(t\) & \(y\) & \(q\) & Weight condition \\
\midrule\noalign{}
\endhead
\bottomrule\noalign{}
\endlastfoot
1 & 3 & 9 & \(S\ge4\) \\
1 & 4 & 8--10 & \(S\ge3\) \\
1 & 5 & 7--11 & \(S\ge2\) \\
2 & 10 & 8 & \(S\ge3\) \\
2 & 11 & 7--9 & \(S\ge2\) \\
3 & 17 & 7 & \(S\ge2\) \\
\end{longtable}

The omitted top offset \(y=(2g-2)t\) gives a
nonpositive-parabolic-degree dual quotient of the stable adjoint bundle
and is impossible for that reason.

\section{Genus-four branch analysis}\label{genus-four-branch-analysis}

\subsection{The first no-high branch is
eliminated}\label{the-first-no-high-branch-is-eliminated}

In the \((S,q)=(0,7)\) branch, \(G\) has rank seven, degree one, and
four sections. If \(H\subsetneq G\) is a subbundle and
\(\widetilde H\subset E\) its inverse image, stability of the
degree-zero bundle \(E\) gives \(\deg\widetilde H<0\), hence
\(\deg H\le0\). Thus \(G\) is stable. Its saturated evaluation image is
generically generated of degree zero, so it is \(W=\mathcal O_C^4\). The
quotient is stable as well: a proper subbundle of \(Q=G/W\) pulls back
to a subbundle of \(G\) of degree at most zero. Therefore

\[L=M^{-1},\qquad
W=\mathcal O_C^4\subset G,\qquad
(\mathrm{rk}\,Q,\deg Q)=(3,1).\]

Let \(P\subset E\) be the inverse image of \(W\). Since \(H^0(P)=0\) and
\(h^1(M)=4\), the boundary of

\[0\longrightarrow M\longrightarrow P\longrightarrow W\longrightarrow0\]

is an isomorphism \(e:H^0(W)\simeq H^1(M)\). Define the quotient-bracket
component

\[f:M\otimes W\longrightarrow Q,\qquad
\phi:H^0(W)\longrightarrow H^0(Q\otimes L).\]

The map \(b:\Lambda^2P\to Q\) supplies information absent from pointwise
Lie algebra. A \v{C}ech comparison of local splittings gives

\[f(v)_*e(u)=f(u)_*e(v)\quad\text{in }H^1(Q).\]

After Serre duality and \(H^0(L\omega_C)\simeq W^\vee\), put
\(V_\phi=W/\ker\phi\) and \(k=\dim V_\phi\). The induced map

\[\Theta:H^0(Q^\vee\omega_C)
\longrightarrow H^0(L\omega_C)\otimes V_\phi^\vee\]

has image in \(\mathrm{Sym}^2(V_\phi^\vee)\).

This symmetry forces \(k\le1\). Here is the complete rank analysis. Let
\(I\subset Q\) be the saturated image of \(f\), of rank \(r\).

\begin{itemize}
\item
  If \(r=1\), stability gives \(\deg I\le0\), so the line \(I\otimes L\)
  has at most one section.
\item
  If \(r=2\), then \(-2\le\deg I\le0\). The quotient line \(R=Q/I\)
  satisfies \(h^0(R^\vee\omega_C)\le3\), while
  \(h^0(Q^\vee\omega_C)\ge8\). Thus \(\mathrm{rk}\,\Theta\ge5\),
  excluding \(k=2\). If \(k\ge3\), three generically spanning sections
  of the rank-two bundle \(I\otimes L\) have at least two independent
  wedges. Otherwise their Pl\"ucker kernel would be constant and give a
  global linear relation. This would be a pencil on a line of degree at
  most two, impossible on the nonhyperelliptic genus-four fibre.
\item
  If \(r=3\) and \(k=3\), then \(\Theta\) is injective, but its domain
  has dimension at least eight and
  \(\dim\mathrm{Sym}^2(\mathbb C^3)=6\).
\item
  If \(r=3\) and \(k=4\), write \(H=Q\otimes L\), let
  \(a:\mathcal O_C^4\to H\), set \(A=\ker a\), and let \(D\) be the
  Fitting divisor of its torsion cokernel, of length \(\ell\). Then
  \(\deg A=\ell-4\) and injectivity on constant sections gives
  \(\ell\le3\). For the raw image \(I_0=\mathrm{im}(a)\), local Smith
  form gives \(\mathcal O_C(-D)I_0^\vee\subset H^\vee\). If
  \(K_A=\ker(W^\vee\to H^0(A^\vee))\), the lower dimension bounds for
  \(H^0(L\omega_C(-D))\) and \(K_A\) are respectively

  \[(4,3,2,1)\quad\text{and}\quad(1,2,2,3)\]

  for \(\ell=0,1,2,3\). Choose nonproportional elements \(u\) and
  \(\eta\) in these spaces. Their product is a section
  \(\alpha\in H^0(Q^\vee\omega_C)\) with
  \(\Theta(\alpha)=u\otimes\eta\). A pure tensor is symmetric only when
  its factors are proportional, a contradiction.
\end{itemize}

Consequently \(\mathrm{rk}\,\phi\le1\). Put \(W_0=\ker\phi\), so
\(\dim W_0\ge3\). If \(h^0(L)=0\), the remaining self-block of \(\beta\)
on \(W_0\) vanishes. If \(L\simeq\mathcal O_C(p)\), write that block as
\(sA\), where \(s\) is the unique section of \(L\) and \(A:W_0\to W\) is
constant. Compatibility with the universal extension gives

\[W_0\xrightarrow{A}W\xrightarrow{e}H^1(M)
\xrightarrow{s_*}H^1(\mathcal O_C),\]

with zero composite. The last map is an isomorphism by the exact
sequence for \(0\to\mathcal O_C(-p)\to\mathcal O_C\to k(p)\to0\), and
\(e\) is an isomorphism, so \(A=0\). Thus in either case
\(\mathcal O_C^3\subset\ker\beta\).

There is a second, shorter global obstruction once three kernel sections
are available. Put \(K=N_E(M)/M\). The normalizer acts on \(M\) through
a character \(\lambda:K\to\mathcal O_C\). For global \(x,y\in H^0(K)\),
local lifts satisfy

\[\partial[x,y]=\lambda(x)\partial y-\lambda(y)\partial x.\]

Here \(\partial:H^0(G)\to H^1(M)\) is injective because \(H^0(E)=0\).
Hence

\[[x,y]=\lambda(x)y-\lambda(y)x.\]

No generically independent three-dimensional subspace of an
\(\mathfrak{sl}_3\) quotient normalizer obeys this law. Regular orbit
types have quotient normalizer dimension at most two. In the semisimple
\((2,1)\) case the quotient is \(\mathfrak{sl}_2\) and \(\lambda=0\), so
the law would make it abelian. For a minimal nilpotent \(m=E_{12}\), use
the quotient basis

\[D=\mathrm{diag}(1/2,-1/2,0),\quad
H=\mathrm{diag}(1/2,1/2,-1),\quad E_{13},\quad E_{32}.\]

The character is one on \(D\) and zero on the other three vectors. Its
kernel is nonabelian, and its only two-dimensional abelian subalgebra is
\(\langle E_{13},E_{32}\rangle\). If a three-space is not contained in
that kernel, the displayed law requires eigenvalue one with multiplicity
two on this two-dimensional intersection. The actual eigenvalues are
\((1+3c)/2\) and \((1-3c)/2\) for \(D+cH+aE_{13}+bE_{32}\), and cannot
both be one. This contradiction eliminates \((S,q)=(0,7)\) without a
determinant or relative-section endpoint.

\subsection{\texorpdfstring{The next no-high branch:
\((S,q)=(0,8)\)}{The next no-high branch: (S,q)=(0,8)}}\label{the-next-no-high-branch-sq08}

Now \((\mathrm{rk}\,G,\deg G,h^0(G))=(7,2,5)\) and \(\deg L=2\).
Positive HN degrees give exactly

\begin{longtable}[]{@{}lr@{}}
\toprule\noalign{}
HN type & Section maximum \\
\midrule\noalign{}
\endhead
\bottomrule\noalign{}
\endlastfoot
\((7,2)\) & 5 \\
\((1,1)+(6,1)\) & \(1+4\) \\
\((2,1)+(5,1)\) & \(1+4\) \\
\((3,1)+(4,1)\) & \(2+3\) \\
\end{longtable}

The evaluation saturation is consequently either

\[\mathcal O_C^5\subset G,\qquad
(\mathrm{rk}\,Q,\deg Q)=(2,2),\]

with \(Q\) semistable, or

\[J\oplus\mathcal O_C^4\subset G,\qquad
\deg J=1,\qquad
(\mathrm{rk}\,R,\deg R)=(2,1),\]

with \(R\) stable. In the second case the \v{C}ech argument must use the
full evaluation injection \(\mathcal O_C^5\to G\), including the section
of \(J\); using only \(\mathcal O_C^4\) loses the boundary isomorphism.

We use two elementary linear-algebra lemmas to make the remaining
argument explicit.

First, let \(V\subset H^0(H)\) have dimension \(k\), where \(H\) has
rank two, and consider

\[\delta:\Lambda^2V\longrightarrow H^0(\det H).\]

If \(\ker\delta\) contains a decomposable vector \(s\wedge t\), the
saturation of \(s,t\) is a line subbundle with at least two sections.
Otherwise \(\mathbb P(\ker\delta)\) is disjoint from
\(\mathrm{Gr}(2,V)\subset\mathbb P(\Lambda^2V)\). Since the Grassmannian
has dimension \(2k-4\), the projective dimension theorem gives

\[\dim\mathrm{im}(\delta)\ge2k-3.\]

Second, a projective pencil of singular ternary quadrics has either a
common linear factor or a common kernel point. Indeed, a nontrivial such
pencil has a rank-two member: otherwise two independent rank-one members
would have a rank-two sum. Such a member can be written \(xy\). Write
the other member as

\[ax^2+by^2+cz^2+dxy+exz+fyz.\]

The determinant of the pencil vanishes identically only if \(c=0\),
\(ef=0\), and \(af^2+be^2=0\). If \(e=f=0\), the pencil has the common
kernel point \([0:0:1]\); if only \(e\) or only \(f\) vanishes, it has
the common factor \(y\) or \(x\), respectively.

Return first to the \(\mathcal O_C^5\) evaluation case. Put
\(H=Q\otimes L\), let \(V=\mathrm{im}\,\phi\), and set \(k=\dim V\). The
same zero-weight condition gives \(\det E\simeq\mathcal O_C\), hence
\(\det G\simeq L\) and \(\det Q\simeq L\). The same symmetric-\(\Theta\)
argument gives the initial dichotomy: a rank-one saturated bracket image
has \(k\le2\), with equality only when \(I\otimes L\) is a degree-three
trigonal pencil; a rank-two image cannot have \(k=2\), so it has
\(k\ge3\).

There is one degree-two issue before \(\ker\phi\) becomes a normalizer
kernel. If \(W_0=\ker\phi\) and \(L\) is ineffective, the remaining
self-block vanishes and \(W_0\otimes\mathcal O_C\subset\ker\beta\). If
\(L\simeq\mathcal O_C(D)\) is effective, its unique section gives a
self-block \(sA_0\). Extension compatibility only forces the image of
\(A_0\) into the one-dimensional kernel of

\[s_*:H^1(M)\longrightarrow H^1(\mathcal O_C),\]

so at least \(\dim W_0-1\) normalizer sections are guaranteed. The full
boundary is injective because \(H^0(E)=0\). Thus every case with
\(k\le1\) supplies at least three genuine normalizer sections and is
eliminated by the boundary-cocycle obstruction. If \(k=2\), the
ineffective-\(L\) case is eliminated as well, leaving only an apparently
exceptional effective-\(L\) case. That exception is also impossible, but
its lift defect must be retained.

Let \(I\subset Q\) be the degree-one bracket image, put
\(A=I\otimes L\), and let \(R_1=Q/I\). Then \(A\) is a basepoint-free
\(g^1_3\). Let

\[T=\mathrm{im}\bigl(H^0(Q^\vee\omega_C)
  \longrightarrow H^0(I^\vee\omega_C)\bigr).\]

For nonzero \(\alpha\in T\), symmetry says that the two-space
\(\alpha H^0(A)\) equals the fixed plane
\(V^\vee\subset H^0(L\omega_C)\). Two independent elements of \(T\)
would therefore have a ratio preserving \(H^0(A)\). Multiplication by
that ratio has an eigenvector, so the ratio is constant on the
nonvanishing locus of the eigenvector and hence globally. Thus
\(\dim T\le1\). Exactness and Riemann--Roch give

\[\dim T=(4+h^0(Q))-(2+h^0(R_1))
      =2+h^0(Q)-h^0(R_1).\]

Since \(h^0(R_1)\le1\), one must have \(R_1\simeq\mathcal O_C(r)\) and
\(h^0(Q)=0\). The inclusion \(I\subset Q\) then gives \(h^0(I)=0\).
Determinants give \(I\otimes R_1=L\), hence \(I=L(-r)\) and
\(A=L^2(-r)\). The equality characterization in
Clifford's theorem gives \(h^0(L^2)\le2\) on a
nonhyperelliptic genus-four curve. Since \(h^0(A)=2\), the point \(r\)
would be a base point of \(|L^2|\). But \(h^0(I)=0\) says that \(r\) is
not in the divisor of the unique section of the effective degree-two
line \(L\), so the square of that section does not vanish at \(r\). This
contradiction eliminates the rank-one \(k=2\) case.

It remains to exclude full bracket image in \(Q\). For \(k=3\),
\(\Theta\) is injective and

\[\dim\mathrm{im}\,\Theta=4+h^0(Q),
\qquad h^0(Q)\le2.\]

On the dense open where \(V\otimes\mathcal O_C\to H\) has rank two, let
\(\ell_x\in\mathbb P(V)\) be its kernel and let \(u_x\in\mathbb P(V)\)
be the evaluation vector of the fixed inclusion
\(V^\vee\subset H^0(L\omega_C)\). If \(M_\alpha\) is the symmetric
matrix \(\Theta(\alpha)\), evaluating on \(\ell_x\) gives

\[u_x^tM_\alpha\ell_x=0.\]

Consequently the singular conic \(q_x=\mathrm{sym}(u_x\otimes\ell_x)\)
lies in \((\mathrm{im}\,\Theta)^\perp\subset\mathrm{Sym}^2(V)\). If
\(h^0(Q)=2\) this orthogonal space is zero. If \(h^0(Q)=1\) it is a
fixed projective point, and unique factorization of the fixed reducible
conic makes either \([u_x]\) or \([\ell_x]\) constant, contradicting
independence of the dual sections or of \(V\). If \(h^0(Q)=0\), the
orthogonal space is a projective pencil. The map \(x\mapsto[q_x]\)
cannot be constant by the same factorization argument, so the whole
pencil is singular. The second lemma gives either a common factor, again
making one of \([u_x],[\ell_x]\) constant, or a common kernel point. In
the latter case all \(u_x\) lie in a fixed two-plane, so a nonzero
member of \(V^\vee\) vanishes identically. Every alternative is
impossible.

For \(k=4\) or \(5\), apply the wedge lemma to \(H\). Here \(H\) is
semistable of rank two and degree six, and \(\det H=L^3\) has at most
four sections. If \(k=5\), the lemma forces a line subbundle
\(A_1\subset H\) with at least two sections. Semistability and
nonhyperellipticity force \(\deg A_1=3\) and \(h^0(A_1)=2\); the
degree-three quotient has at most two sections, contradicting \(k=5\).

If \(k=4\), the same lemma forces a line subbundle, and equality in the
section bounds forces an exact sequence

\[0\longrightarrow A_1\longrightarrow H\longrightarrow A_2
\longrightarrow0\]

in which \(A_1,A_2\) are basepoint-free \(g^1_3\)'s and
every section of \(A_2\) lifts. A base point in either pencil would
leave a degree-two pencil and make \(C\) hyperelliptic. In particular
\(V=H^0(H)=V_1\oplus V_2\), with \(V_1=H^0(A_1)\) and \(V_2\) a lift of
\(H^0(A_2)\). Twisting by \(M\) gives \(I=A_1\otimes M\subset Q\) and
\(R_2=A_2\otimes M=Q/I\). For
\(\alpha\in H^0(R_2^\vee\omega_C)\subset H^0(Q^\vee\omega_C)\), the
\(V_1\) columns of \(\Theta(\alpha)\) vanish; symmetry makes its \(V_1\)
rows vanish as well. Hence \(\alpha H^0(A_2)\) is the fixed two-plane
\(V_2^\vee\). But

\[h^0(R_2^\vee\omega_C)=2+h^0(R_2)\ge2,\]

so two independent choices of \(\alpha\) give the same product pencil.
Their ratio preserves \(H^0(A_2)\) and is constant by the eigenvector
argument, a contradiction. Thus every \(\mathcal O_C^5\) evaluation case
is eliminated.

For the \(J\oplus\mathcal O_C^4\) evaluation, the full map
\(\mathcal O_C^5\to G\) is not saturated at the zero of the section of
\(J\). It is enough for the \v{C}ech symmetry, but bracket with \(M\) is not
known to preserve that raw evaluation lattice. Its self-block may
acquire a pole in a raw trivialization. The full-rank residual cases can
nevertheless be excluded before confronting that lattice issue. Write
the stable residual bundle as \(R_0\) and put \(H_0=R_0\otimes L\), a
stable rank-two bundle of degree five. If \(H_0\) had four sections, the
wedge lemma and \(h^0(\det H_0)\le3\) would produce a line subbundle
with two sections; stability gives degree at most two, contradicting
nonhyperelliptic gonality. Thus \(h^0(H_0)\le3\).

A full bracket image cannot have coefficient rank two, because
injectivity of \(\Theta\) gives

\[\dim H^0(R_0^\vee\omega_C)=5+h^0(R_0)>3
  =\dim\mathrm{Sym}^2(\mathbb C^2).\]

It must therefore have \(k=3\). The same \(q_x\) construction now has
\(\dim\mathrm{im}\,\Theta=5\) or \(6\): indeed, stability gives
\(h^0(R_0)\le1\). To see this, two sections with zero wedge would lie in
a degree-zero line with two sections, while a nonzero wedge vanishes at
the unique point of \(\det R_0\) and a linear combination would generate
a degree-at-least-one line subbundle. Both alternatives contradict
stability. Thus the orthogonal of \(\mathrm{im}\,\Theta\) in
\(\mathrm{Sym}^2(\mathbb C^3)\) is projectively a point or empty. The
empty case is immediate; in the point case unique factorization again
forces an evaluation vector or a kernel line to be constant. Hence every
full-rank residual image is impossible. A rank-one residual image has
\(k\le1\): its saturated line in \(R_0\) has degree at most zero, so
after twisting by \(L\) it has degree at most two and at most one
section. It remains to convert the resulting large kernel into genuine
normalizer sections without pretending that the raw evaluation lattice
is saturated.

Let

\[V_0=H^0(G),\qquad
F=\operatorname{im}(V_0\otimes\mathcal O_C\to G),\qquad
W=J\oplus\mathcal O_C^4.\]

The five sections are generically independent, so
\(F\simeq\mathcal O_C^5\), and the degree-one saturation defect gives
\(W/F\simeq k(p)\). Stability gives \(H^0(E)=0\); since \(h^1(M)=5\),
the extension boundary is an isomorphism

\[e:V_0\xrightarrow{\sim}H^1(M).\]

Write \(\phi:V_0\to H^0(R_0\otimes L)\) for the residual
quotient-bracket map. The \v{C}ech symmetry is valid for the actual raw
sections: choosing local lifts in \(E\) and projecting their bracket to
\(R_0\) gives

\[\phi(v)_*e(u)=\phi(u)_*e(v)
\quad\text{in }H^1(R_0).\]

Suppose first that \(\dim\operatorname{im}\phi=1\). Write
\(\phi(v)=c(v)r\) with \(0\ne r:M\to R_0\). Taking \(u\in\ker c\) and
\(v\) with \(c(v)=1\) shows that \(r_*e(u)=0\), hence

\[\dim\ker\bigl(H^1(M)\xrightarrow{r_*}H^1(R_0)\bigr)\ge4.\]

Let \(I\subset R_0\) be the saturation of \(r(M)\) and put
\(d=\operatorname{length}(I/r(M))\). Stability of the rank-two,
degree-one bundle \(R_0\) gives \(\deg I\le0\), so \(0\le d\le2\). For
\(Q_r=R_0/r(M)\) there is an exact sequence

\[0\longrightarrow I/r(M)\longrightarrow Q_r
\longrightarrow R_0/I\longrightarrow0,
\qquad \deg(R_0/I)=3-d.\]

On the nonhyperelliptic genus-four test fibre, \(h^0(R_0/I)\le2,1,1\)
for \(d=0,1,2\), respectively. Thus \(h^0(Q_r)\le2,2,3\). The long exact
sequence of \(0\to M\to R_0\to Q_r\to0\) gives

\[\dim\ker r_*=h^0(Q_r)-h^0(R_0)\le h^0(Q_r)\le3,\]

contradicting the preceding lower bound. Therefore \(\phi=0\).

Now \(\beta(V_0)\subset H^0(W\otimes L)\). Compose with the
principal-part map to \(H^0((W/F)\otimes L)\) and let \(V_1\) be its
kernel. The target is one-dimensional, so \(\dim V_1\ge4\), and

\[\beta(V_1)\subset H^0(F\otimes L)
=V_0\otimes H^0(L).\]

If \(h^0(L)=0\), all of \(V_1\) consists of genuine normalizer sections.
If \(h^0(L)=1\), write its generator as \(s\) and \(\beta|_{V_1}=sA\)
for a linear map \(A:V_1\to V_0\). The global bracket map
\(M\otimes G\to E\) lifts \(sA(x):M\to G\) for every \(x\in V_1\). Hence
the pullback of the extension along multiplication by \(s\) splits, or
\(s_*e(Ax)=0\).

The zero divisor of \(s\) has length two, and
\(0\to M\xrightarrow{s}\mathcal O_C\to\mathcal O_D\to0\) gives
\(\dim\ker(s_*:H^1(M)\to H^1(\mathcal O_C))=1\). Since \(e\) is an
isomorphism, \(\operatorname{rk}A\le1\) and therefore
\(\dim\ker A\ge3\). These are three generically independent global
normalizer sections because \(F\simeq V_0\otimes\mathcal O_C\). The
normalizer boundary-cocycle obstruction excludes them. Thus the entire
nonsaturated evaluation locus, and hence all of \((S,q)=(0,8)\), is
impossible.

\subsection{\texorpdfstring{The zero-weight \(q=9\) branch and its
closure}{The zero-weight q=9 branch and its closure}}\label{the-zero-weight-q9-branch-and-its-closure}

For \((S,q)=(0,9)\) one has

\[\deg M=-3,\qquad \deg L=3,\qquad
(\operatorname{rk}G,\deg G,h^0(G))=(7,3,6).\]

Stability gives \(H^0(E)=0\), so the full boundary is an isomorphism

\[e:H^0(G)\xrightarrow{\sim}H^1(M)\]

between six-dimensional spaces. The exact HN equality types allowed by
the small-slope bounds are

\[(7,3),\qquad (1,1)+(6,2),\qquad (2,2)+(5,1).\]

The saturation \(W\) of the six-section evaluation has rank six.
Stability of \(E\) bounds its degree by two, and equality in the section
bounds gives exactly

\[\mathcal O_C^6,\qquad
J_1\oplus\mathcal O_C^5,\qquad
A_2\oplus\mathcal O_C^5,\qquad
B_{2,2}\oplus\mathcal O_C^4.\]

Here \(J_1\) is a degree-one line with one section, \(A_2\) is a
degree-two line with one section, and \(B_{2,2}\) is a semistable
rank-two, degree-two bundle with two sections. If
\(d=\deg W\in\{0,1,2\}\) and
\(F=\operatorname{im}(H^0(G)\otimes\mathcal O_C\to G)\), then

\[F\simeq\mathcal O_C^6,\qquad
\operatorname{length}(W/F)=d,
\qquad R:=G/W\simeq L(-D),\quad \deg D=d.\]

Let

\[\phi:H^0(G)\longrightarrow H^0(R\otimes L),
\qquad U=\operatorname{im}\phi,\quad k=\dim U\]

be the residual bracket coefficient. If \(\phi\ne0\), \v{C}ech symmetry
gives an injective map

\[\Theta:H^0(R^{-1}\otimes\omega_C)
\longrightarrow\operatorname{Sym}^2(U^\vee).\]

Every nonzero form in its image is nondegenerate. Indeed, choose
\(0\ne u\in U\) and a lift \(v\in H^0(G)\) with \(\phi(v)=u\). If \(u\)
were a null vector, then for every \(w\in H^0(G)\) the \v{C}ech symmetry
would give

\[0=\langle e(v),\alpha\phi(w)\rangle
 =\langle e(w),\alpha u\rangle.\]

The boundary isomorphism and Serre duality would force \(\alpha u=0\),
which is impossible for two nonzero sections on an integral curve. A
complex pencil of symmetric forms always contains a singular member, so

\[\phi\ne0\quad\Longrightarrow\quad
h^0(R^{-1}\otimes\omega_C)\le1.\]

Riemann--Roch and \(R=L(-D)\) give the exact dimension

\[h^0(R^{-1}\otimes\omega_C)=d+h^0(R).\]

Consequently \(d=2\) forces \(k=0\); \(d=1\) forces \(k=0\) when \(R\)
is effective; and \(d=0\) forces \(k=0\) when \(L\) is trigonal.

The remaining issue is no longer the residual Petri map. After taking
\(\ker\phi\) and killing the \(d\) principal parts of \(W/F\), there is
a subspace \(V_1\subset H^0(G)\) of dimension at least \(6-k-d\) on
which the self-block is

\[V_1\longrightarrow H^0(G)\otimes H^0(L).\]

Write \(\ell=h^0(L)\). On the nonhyperelliptic genus-four test fibre,
\(\ell\in\{0,1,2\}\). Extension compatibility places the self-block in
the kernel of

\[\delta_L:H^1(M)\otimes H^0(L)\longrightarrow H^1(\mathcal O_C),
\qquad a\otimes s\longmapsto s_*a.\]

For \(\ell=0,1,2\), the kernel dimensions are respectively

\[\kappa_\ell=0,2,8.\]

For \(\ell>0\) this follows because multiplication by any nonzero
section of \(L\) surjects \(H^1(M)\to H^1(\mathcal O_C)\); the domain of
\(\delta_L\) has dimension \(6\ell\) and the target has dimension four.
The guaranteed normalizer rank is therefore only

\[N=\max\{0,6-k-d-\kappa_\ell\}.\]

This formula closes some subloci, but not the branch:

\begin{longtable}[]{@{}lrrl@{}}
\toprule\noalign{}
Evaluation data & Residual rank & Guaranteed \(N\) & Status from the
extension count alone \\
\midrule\noalign{}
\endhead
\bottomrule\noalign{}
\endlastfoot
\(d=2\), \(\ell=0\) & \(k=0\) & \(4\) & closed \\
\(d=2\), \(\ell=1\) & \(k=0\) & \(2\) & not closed at this stage \\
\(d=2\), \(\ell=2\) & \(k=0\) & \(0\) & not closed at this stage \\
\(d=1\), \(R\) effective, \(\ell=1\) & \(k=0\) & \(3\) & closed \\
\(d=1\), \(R\) effective, \(\ell=2\) & \(k=0\) & \(0\) & not closed at
this stage \\
\(d=0\), \(\ell=2\) & \(k=0\) & \(0\) & not closed at this stage \\
\end{longtable}

There are also possible escapes with nonzero residual coefficient rank:
\(d=0,\ell=0,k=4\) when \(h^0(L^2)=4\), equivalently
\(L^2\simeq\omega_C\); \(d=0,\ell=1,k\ge2\); and, when \(R\) is
ineffective, \(d=1,\ell=0,k=3\) when \(h^0(R\otimes L)=3\), or
\(d=1,\ell=1,k\ge1\). The case \(d=1\), \(R\) ineffective, \(\ell=2\)
cannot occur because a degree-three pencil on the nonhyperelliptic fibre
is basepoint-free. In each surviving case the current lower bound can
fall below three.

The global Killing pairing supplies the missing constraint. Let
\(Q_2=\operatorname{tr}(m^2)\in H^0(L^2)\). Pairing the inverse image
\(P\subset E\) with \(M\) by the global Killing form gives

\[b:P\longrightarrow M^*=L,
\qquad b|_M=Q_2.\]

Thus the pushout of the universal extension
\(0\to M\to P\to\mathcal O_C^6\to0\) along \(Q_2:M\to L\) splits. Since
its boundary \(e:\mathbb C^6\to H^1(M)\) is an isomorphism,
multiplication

\[(Q_2)_*:H^1(M)\longrightarrow H^1(L)\]

must vanish. If \(Q_2\ne0\), its Serre-dual map is multiplication by \(Q_2\)
from \(H^0(\omega_C\otimes L^{-1})\) to \(H^0(\omega_C\otimes L)\); this
is injective on the integral curve and has nonzero source, of dimension
\(h^0(L)\ge1\). Hence \((Q_2)_*\) cannot vanish. Therefore every
effective-\(L\) realization,
equivalently every sublocus with \(h^0(L)>0\), must satisfy

\[Q_2=\operatorname{tr}(m^2)=0.\]

In the unpunctured no-high branch, Killing gives more than isotropy.
Since \(Q_2=0\), pairing with \(M\) descends to

\[\bar b:G=E/M\longrightarrow L.\]

For \(x,y\in H^0(G)\), choose local lifts with boundary cocycles
representing \(e(x),e(y)\in H^1(M)\). The overlap formula for their
Killing pairing gives

\[\bar b(y)_*e(x)+\bar b(x)_*e(y)=0
\quad\text{in }H^1(\mathcal O_C).\]

If the induced map \(H^0(G)\to H^0(L)\) were nonzero, choose \(x\) whose
image is a nonzero section \(s\). Its kernel has dimension at least
\(6-h^0(L)\ge4\), and the displayed identity would carry that kernel
injectively under \(e\) into

\[\ker\bigl(s_*:H^1(M)\longrightarrow H^1(\mathcal O_C)\bigr),\]

which has dimension two. This is impossible. Hence \(\bar b\) kills the
raw evaluation \(F\simeq\mathcal O_C^6\). If \(P\subset E\) is its
inverse image, then \(P\subset M^\perp\); both bundles have rank seven
and degree \(-3\), so

\[P=M^\perp,
\qquad F=M^\perp/M.\]

In particular \(F\) is saturated. Thus every effective-\(L\) case with
\(d=1\) or \(d=2\) is impossible. Moreover
\([M,M^\perp]\subset M^\perp\), so the residual coefficient map vanishes
automatically when \(d=0\). If \(h^0(L)=1\), the remaining self-block is
\(sA\) and compatibility puts \(e(\operatorname{im}A)\) in the
two-dimensional kernel of \(s_*\). Hence \(\operatorname{rk}A\le2\) and
at least four constant normalizer sections remain, a contradiction. The
last effective case left by these extension counts is

\[d=0,\qquad h^0(L)=2,\qquad P=M^\perp.\]

Here Killing descends to a nondegenerate constant symmetric form on
\(F\simeq\mathcal O_C^6\), and the self-block is a pencil skew for that
form.

The actual inner \(\mathrm{PGL}_3\) structure now gives a topological
contradiction. The adjoint homomorphism

\[\mathrm{PGL}_3\longrightarrow\mathrm{SO}_8\]

lifts to \(\mathrm{Spin}_8\): lift first from the simply connected group
\(\mathrm{SL}_3\), and observe that its central \(\mu_3\) must map
trivially to the \(\mu_2\) kernel of
\(\mathrm{Spin}_8\to\mathrm{SO}_8\). Hence the orthogonal obstruction
class of the adjoint bundle is \(w_2(E)=0\). On the other hand, a smooth
orthogonal splitting along the isotropic line gives

\[E\simeq H(M)\perp F,
\qquad
w_2(E)=c_1(M)+w_2(F)\pmod2.\]

The form on the trivial bundle \(F\simeq\mathcal O_C^6\) is a constant
nondegenerate matrix, so \(w_2(F)=0\), while \(\deg M=-3\) is odd. This
gives \(w_2(E)=1\), a contradiction. Here \(w_2\) is the obstruction
class of the complex orthogonal bundle, not the Stiefel--Whitney class
of its underlying realification. Therefore \textbf{every effective-\(L\)
no-high \(q=9\) sublocus is impossible}.

The punctured canonical lattices admit an even sharper elementary
obstruction. At a marked DVR their standard basis consists of two Cartan
generators and root generators

\[z^{\epsilon_{ij}}E_{ij},
\qquad \epsilon_{ij}\in\{0,1\}.\]

For an opposite root pair, \(\epsilon_{ij}+\epsilon_{ji}\) is zero
inside a Levi block and one across distinct flag blocks. Let a
saturated, hence primitive, isotropic generator \(m\) have pairing ideal
\(\operatorname{tr}(mE)=z^{d_p}R\). If a unit coefficient of \(m\)
occurs in the Cartan part, the unimodular Cartan Gram matrix gives
\(d_p=0\). If it occurs on a root generator, pairing with the unique
opposite root generator isolates that coefficient and gives \(d_p\le1\).
Thus

\[d_p\le1\]

at every nonscalar marking, with no possible cancellation; at scalar or
ordinary points \(d_p=0\).

Let \(N_{\rm ns}\) be the number of nonscalar markings and let \(s_p\)
be the adjoint determinant contribution, two for type \(2+1\) and three
for a full flag. If \(D_b=\sum_p d_p p\), then

\[\deg D_b\le N_{\rm ns}\le S/2,
\qquad S=\sum s_p.\]

On the other hand,

\[\operatorname{im}(E\longrightarrow M^*)=L(-D_b),
\qquad
\deg M^\perp=\deg D_b-S-3.\]

The symmetric \v{C}ech identity gives \(P\subset M^\perp\), while
\(\deg P=-3\), and therefore \(\deg D_b\ge S\). This contradicts the
displayed upper bound whenever \(S>0\). Hence \textbf{every effective
punctured high-HN \(q=9\) row is impossible}. This argument uses the
canonical periodic-chain root exponents; it is not asserted for an
arbitrary commensurable lattice.

The ineffective-\(L\) subloci close by a different two-normalizer
argument. Put \(\ell=h^0(L)=0\). The extension count already closes
\(d=2\). For \(d=1\), the only possible escape has \(k=3\): the
three-dimensional space \(\ker\phi\) loses at most one dimension to the
length-one principal-part map. Its remaining two-space has bracket in
\(H^0(F\otimes L)=H^0(L)^6=0\). For \(d=0\), the only escape has \(k=4\)
and the two-dimensional space \(\ker\phi\) has the same property. Thus
in either case there are two generically independent global sections

\[\mathcal O_C^2\longrightarrow K:=\ker\beta.\]

The two maximal escapes have the exact line-bundle
forms

\[(d,k)=(0,4):\quad L^2\simeq\omega_C,\]

and

\[(d,k)=(1,3):\quad
L^2\simeq\omega_C(D-p)\]

for an effective divisor \(D\) of degree one and a point \(p\). Lower
coefficient rank already leaves at least three normalizer sections.

Let \(Q_2=\operatorname{tr}(m^2)\in H^0(L^2)\). If \(Q_2=0\), Killing
pairing descends on the inverse image \(P\) of the raw
\(F\simeq\mathcal O_C^6\) to a map \(F\to L\). This map is zero because
\(h^0(L)=0\), so \(P\subset M^\perp\). Equal rank and degree give
\(P=M^\perp\). For \(d=1\) this already contradicts the nonsaturation of
\(F\); for \(d=0\) it gives \(M^\perp/M\simeq\mathcal O_C^6\) and the
Spin-parity contradiction above.

It remains to take \(Q_2\ne0\). For a quotient-normalizer element \(x\),
write \([x,m]=\lambda(x)m\). Killing invariance gives

\[\lambda(x)Q_2=\operatorname{tr}([x,m]m)=0,\]

so the grading character \(\lambda\) vanishes. The injective-boundary
\v{C}ech law then makes the two global normalizer sections commute. After
extending the generic field if necessary, if \(m\) is regular, its
centralizer has dimension two and \(Z(m)/\mathbb C m\) has dimension
one, so it cannot contain them. If \(m\) is nonregular and \(Q_2\ne0\),
its minimal polynomial has degree at most two and tracelessness forces
the semisimple eigenvalue pattern \((a,a,-2a)\) with \(a\ne0\). In that
case

\[Z(m)/\mathbb C m\simeq\mathfrak{sl}_2,\]

which has no two-dimensional abelian subalgebra. This is the final
contradiction. Therefore \textbf{the entire unpunctured no-high branch
\((S,q)=(0,9)\) is impossible}.

\subsection{\texorpdfstring{The first punctured branch:
\((S,q)=(2,8)\)}{The first punctured branch: (S,q)=(2,8)}}\label{the-first-punctured-branch-sq28}

Here the branch data are

\[\deg M=-2,\qquad \deg L=2,\qquad
(\mathrm{rk}\,G,\deg G,h^0(G))=(7,0,5),
\qquad h^1(M)=5.\]

The no-high bundle \(A\) has rank seven, degree \(42\), and
\(\mu_{\max}(A)\le6=\mu(A)\). Hence \(A\) is semistable, and
\(G=A^\vee\otimes\omega_C\) is semistable of degree zero. The saturated
evaluation image is generically generated, so it has nonnegative degree;
semistability of \(G\) gives the opposite inequality. It is therefore a
degree-zero trivial bundle. Since all five global sections of \(G\)
factor through it, its rank is at least five; an evaluation image
generated by five sections has rank at most five. Thus its rank is
exactly five:

\[W=\mathcal O_C^5\subset G.\]

The quotient \(Q=G/W\) is semistable of rank two and degree zero.
Parabolic stability gives \(H^0(E)=0\), so the boundary of the
inverse-image extension is an isomorphism

\[e:H^0(W)\xrightarrow{\sim}H^1(M).\]

Put \(H=Q\otimes L\), a semistable rank-two bundle of degree four. On
the nonhyperelliptic test fibre, Riemann--Roch gives

\[h^0(\det H)=1+h^0(\omega_C\otimes\det H^{-1})\le2,\]

because the residual line has degree two. Three independent sections of
\(H\) would give a nonzero kernel for

\[\Lambda^2\mathbb C^3\longrightarrow H^0(\det H).\]

Every nonzero bivector in three variables is decomposable, so two
independent sections would generate a line subbundle of degree at most
two. Such a line cannot have two sections on a nonhyperelliptic curve.
Therefore

\[h^0(Q\otimes L)\le2.\]

Let \(\phi:H^0(W)\to H^0(Q\otimes L)\) be the quotient-bracket
coefficient map and put \(k=\dim\operatorname{im}\phi\). If \(k=2\), the
two image sections generically span \(Q\otimes L\); otherwise they again
generate a forbidden degree-at-most-two pencil. The \v{C}ech symmetry then
makes the Petri map injective with image in the symmetric square:

\[H^0(Q^\vee\otimes\omega_C)
\hookrightarrow\mathrm{Sym}^2(\mathbb C^2)^\vee.\]

Riemann--Roch gives \(h^0(Q^\vee\otimes\omega_C)=6+h^0(Q)\ge6\), while
the target has dimension three. Thus \(k\le1\).

Now \(W_0=\ker\phi\) has dimension at least four. If \(L\) is
ineffective, the remaining self-block vanishes. If
\(L\simeq\mathcal O_C(D)\), write the block as \(sA\). Extension
compatibility gives \(s_*eA=0\), and the exact sequence

\[0\longrightarrow M\xrightarrow{s}\mathcal O_C
\longrightarrow\mathcal O_D\longrightarrow0\]

shows that \(\ker(s_*:H^1(M)\to H^1(\mathcal O_C))\) is one-dimensional.
Since \(e\) is an isomorphism, \(\operatorname{rk}A\le1\). At least
three generically independent constant sections therefore lie in
\(\ker\beta\), and the normalizer boundary-cocycle obstruction
eliminates the branch. No Pfaffian effectivity or relative-section
endpoint is used.

\subsection{\texorpdfstring{The punctured no-high \(q=9\) rows are
eliminated}{The punctured no-high q=9 rows are eliminated}}\label{the-punctured-no-high-q9-rows-are-eliminated}

The two punctured no-high numerical rows at coefficient rank \(q=9\) are
\((S,q)=(2,9),(3,9)\).

For \((S,q)=(2,9)\) one has

\[(\operatorname{rk}G,\deg G,h^0(G))=(7,1,6),
\qquad \deg M=-3,
\qquad h^1(M)=6.\]

Let \(W\) be the saturation of the six-section evaluation and \(F\) its
raw image. The no-high slopes of \(G\) are nonnegative. Since \(G/W\)
also has nonnegative degree, one has \(0\le\deg W\le1\). A
rank-at-most-five bundle of degree at most one that is generically
generated cannot carry six independent sections: after choosing its rank
many independent sections it is either \(\mathcal O_C^r\) or
\(\mathcal O_C(p)\oplus\mathcal O_C^{r-1}\), and in either case has at
most \(r\) sections. Hence \(W\) has rank six. The evaluation
\(H^0(G)\otimes\mathcal O_C\to G\) is therefore generically injective,
so \(F\simeq\mathcal O_C^6\), with the exact alternatives

\[W=F=\mathcal O_C^6
\quad\text{or}\quad
W=\mathcal O_C(p)\oplus\mathcal O_C^5,
\quad W/F\simeq k(p).\]

Write \(\delta=\deg W\in\{0,1\}\) and \(R=G/W\), so \(\deg R=1-\delta\).
If the residual bracket coefficient
\(\phi:\mathbb C^6\to H^0(R\otimes L)\) were nonzero, \v{C}ech symmetry
would inject

\[H^0(R^{-1}\otimes\omega_C)
\longrightarrow\operatorname{Sym}^2(\operatorname{im}\phi)^\vee\]

with every nonzero form in the image nondegenerate. But Riemann--Roch
gives

\[h^0(R^{-1}\otimes\omega_C)
=2+\delta+h^0(R)\ge2,\]

If \(\dim\operatorname{im}\phi=1\), the target of the displayed
injection has dimension one, already contradicting the source dimension.
If \(\dim\operatorname{im}\phi\ge2\), any two-dimensional subspace of
the image is a complex pencil of symmetric forms and contains a singular
member. Thus \(\phi=0\).

If \(L\) is ineffective, killing the at-most-one principal part leaves
at least five raw sections, and their self-block lies in
\(H^0(F\otimes L)=0\). The normalizer boundary law gives a
contradiction. If \(L\) is effective, the quadratic Killing pushout
forces \(Q_2=\operatorname{tr}(m^2)=0\). The symmetric Killing \v{C}ech
identity then gives \(P\subset M^\perp\). Since \(S=2\) consists of one
type-\(2+1\) marking, the canonical local basis gives \(\deg D_b\le1\),
while the global degree comparison gives \(\deg D_b\ge S=2\). This is
impossible.

For \((S,q)=(3,9)\) the data are

\[(\operatorname{rk}G,\deg G,h^0(G))=(7,0,6),
\qquad \deg M=-3,
\qquad h^1(M)=6.\]

Here the no-high inequalities make \(G\) semistable of degree zero, so
its saturated evaluation is

\[W=\mathcal O_C^6,\]

with degree-zero quotient line \(Q\). A nonzero residual coefficient
would again produce a subspace of everywhere nondegenerate symmetric
forms, now of dimension

\[h^0(Q^{-1}\otimes\omega_C)=3+h^0(Q)\ge3,\]

which is impossible. Thus the residual coefficient vanishes. If \(L\) is
ineffective, all six raw sections lie in the quotient normalizer. If
\(L\) is effective, the same Killing argument gives
\(P\subset M^\perp\). The equation \(2a+3b=3\) forces exactly one
full-flag marking, so the canonical local bound is \(\deg D_b\le1\),
while the global inclusion requires \(\deg D_b\ge3\). This is again
impossible.

Therefore \textbf{all no-high genus-four rows at \(q=9\) are
impossible}.

\subsection{\texorpdfstring{The no-high \(q=10\) layer is
eliminated}{The no-high q=10 layer is eliminated}}\label{the-no-high-q10-layer-is-eliminated}

The four rows at \(q=10\) have the common data

\[\deg M=-4,\qquad \deg L=4,\qquad h^1(M)=h^0(G)=7,
\qquad \deg G=4-S.\]

For \(S=0\), the unique HN type attaining the required seven-section
upper bound is

\[(\operatorname{rk},\deg)=(2,3)+(5,1),\]

with section contributions three and four. Let \(H\) be the semistable
rank-two, degree-three first factor. Equality would give \(h^0(H)=3\).
The wedge map

\[\Lambda^2H^0(H)\longrightarrow H^0(\det H)\]

is injective: a nonzero kernel vector is decomposable because
\(\dim H^0(H)=3\), and its two sections would generate a line subbundle
with at least two sections. Such a line has degree at least three on the
nonhyperelliptic test fibre, contradicting semistability of \(H\) of
slope \(3/2\). But \(\det H\) has degree three and at most two sections,
so the three-dimensional exterior square cannot inject. The row
\((0,10)\) is impossible.

Now let \(S=2,3,4\). The exact equality HN types are respectively

\[(1,2)+(6,0)\text{ or }(2,2)+(5,0),
\qquad (1,1)+(6,0),
\qquad (7,0).\]

Equality forces every degree-zero factor to be trivial and the positive
factor to have full-rank evaluation. Hence the seven sections have
generically full-rank raw image

\[F\simeq\mathcal O_C^7\subset G,
\qquad d:=\operatorname{length}(G/F)=4-S=2,1,0.\]

If \(P\subset E\) is the inverse image of \(F\), parabolic stability
gives \(H^0(P)=0\), so the extension boundary is an isomorphism

\[e:\mathbb C^7\overset\sim\longrightarrow H^1(M).\]

Riemann--Roch and nonhyperellipticity give \(h^0(L)\in\{1,2\}\).
Suppose first that \(h^0(L)=1\), with generator \(s\). Killing the \(d\)
principal parts leaves a subspace \(V_1\) of dimension at least \(7-d\)
on which the self-block has the form \(sA\). Extension compatibility
gives

\[s_*eA=0,\]

and the degree-four divisor sequence gives

\[\dim\ker\bigl(s_*:H^1(M)\longrightarrow H^1(\mathcal O_C)\bigr)=3.\]

Thus at least \(7-d-3=4-d\) constant quotient-normalizer directions
remain. For \(S=3,4\) this is at least three, and the normalizer
boundary-cocycle obstruction closes the rows.

For \(S=2\), the count leaves two directions. Put
\(Q_2=\operatorname{tr}(m^2)\). If \(Q_2=0\), the Killing pairing
descends to a map \(G\to L\). If its restriction to the seven raw
sections were nonzero, the symmetric Killing \v{C}ech identity would inject
a subspace of dimension at least six under \(e\) into the
three-dimensional kernel of multiplication by \(s\). Hence that
restriction vanishes. Since \(F\) has full generic rank in \(G\), the
Hodge line would then be orthogonal to the whole generic fibre of \(E\),
contradicting nondegeneracy of the \(\mathfrak{sl}_3\) Killing form.

If instead \(Q_2\ne0\), Killing invariance forces the grading character
to vanish on the two surviving quotient-normalizer sections, and the
boundary law makes them commute. A regular element has
\(\dim Z(m)/\mathbb C m=1\). The only nonregular traceless \(3\times3\)
orbit with \(Q_2\ne0\) is semisimple of type \((2,1)\), where
\(Z(m)/\mathbb C m\simeq\mathfrak{sl}_2\), which contains no abelian
two-plane. This closes the last one-section case.

Finally suppose \(h^0(L)=2\). Riemann--Roch gives
\(h^0(\omega_C\otimes L^{-1})=1\). The Killing map splits the pushout of
the universal extension along \(Q_2:M\to L\), so \((Q_2)_*e=0\). Since
\(e\) is surjective, \((Q_2)_*=0\) on \(H^1(M)\). If \(Q_2\) were
nonzero, its Serre-dual multiplication map from the one-dimensional
space \(H^0(\omega_C\otimes L^{-1})\) would be injective. Therefore
\(Q_2=0\), and the preceding full-rank Killing contradiction applies.
Consequently \textbf{every no-high genus-four row at \(q=10\) is
impossible}.

\subsection{A degree-four factor
lemma}\label{a-degree-four-factor-lemma}

We use the following sharpening repeatedly. If \(H\) is semistable of
degree four and rank \(r=1,2,3\) on the nonhyperelliptic genus-four test
fibre, then

\[h^0(H)\le r+1.\]

For \(r=1\), Riemann--Roch reduces the assertion to the fact that a
degree-two line has at most one section. For \(r=2\), suppose
\(h^0(H)=4\) and put \(V=H^0(H)\). Since \(h^0(\det H)\le2\), the kernel
of

\[\Lambda^2V\longrightarrow H^0(\det H)\]

has dimension at least four. Its projectivization has dimension at least
three and must meet the four-dimensional Grassmannian
\(\operatorname{Gr}(2,4)\subset\mathbb P^5\). A decomposable kernel
vector would give a line subbundle with two sections and hence degree at
least three, contradicting semistability at slope two.

For \(r=3\), suppose \(h^0(H)=5\). Every saturated rank-two subsheaf
\(K\subset H\) has degree at most two and cannot have three sections:
the exterior square of three such sections maps to the
at-most-one-dimensional space \(H^0(\det K)\), so a decomposable kernel
vector would give a two-section line subbundle, impossible under
semistability of \(H\). On the other hand, the kernel of

\[\Lambda^3H^0(H)\longrightarrow H^0(\det H)\]

has dimension at least eight. Its projectivization is at least
\(\mathbb P^7\) and must meet
\(\operatorname{Gr}(3,5)\subset\mathbb P^9\), of dimension six. The
resulting three sections lie in a saturated subsheaf of rank at most
two, which was just excluded. This proves the lemma.

\subsection{\texorpdfstring{The no-high \(q=11\) layer is
eliminated}{The no-high q=11 layer is eliminated}}\label{the-no-high-q11-layer-is-eliminated}

For \((S,q)=(0,11)\) one has

\[(\operatorname{rk}G,\deg G,h^0(G))=(7,5,8).\]

The only HN types whose numerical bounds reach eight sections are

\[(2,4)+(5,1),\qquad (3,4)+(4,1),\]

with respective section bounds \(4+4\) and \(5+3\). The degree-four
factor lemma improves the leading bounds to three and four, so both
totals are at most seven. Thus \((0,11)\) is impossible.

For \((S,q)=(2,11)\) the data are

\[(\operatorname{rk}G,\deg G,h^0(G))=(7,3,8),
\qquad \deg M=-5,\qquad h^1(M)=8.\]

The equality types are

\[(1,3)+(6,0),\qquad (2,3)+(5,0).\]

The rank-two, degree-three factor in the second type cannot have three
sections, by the same exterior-square argument used at \(q=10\).
Equality in the first type forces a splitting

\[G\simeq J\oplus\mathcal O_C^6,\]

where \(J\) is a basepoint-free trigonal pencil. The evaluation sequence
is

\[0\longrightarrow J^{-1}\longrightarrow\mathcal O_C^8
\longrightarrow G\longrightarrow0.\]

The weight \(S=2\) comes from a unique type-\(2+1\) marking \(p\). Since
\(\det E\simeq\mathcal O_C(-2p)\) and \(\det G=J\), one has

\[L=M^{-1}=J(2p).\]

The boundary \(e:H^0(G)\to H^1(M)\) is an isomorphism. Project the
quotient bracket to the \(J\otimes L\) target and call its coefficient
map \(\phi\). \v{C}ech symmetry gives an injection

\[H^0(\omega_C\otimes J^{-1})
\longrightarrow\operatorname{Sym}^2(\operatorname{im}\phi)^\vee\]

whose nonzero image forms are nondegenerate. The source has dimension
two. If \(\dim\operatorname{im}\phi=1\), the target is only
one-dimensional; if \(\dim\operatorname{im}\phi\ge2\), the image
contains a complex pencil and therefore a singular form. Both are
impossible, so \(\phi=0\). Since \(G\) is globally generated, the whole
target-\(J\) bracket block vanishes.

Let \(J'=\omega_C\otimes J^{-1}\) be the residual trigonal pencil.  On a
very general genus-four curve, every trigonal pencil is induced by one of
the two rulings of the smooth canonical quadric: a trigonal pencil produces
a ruling of the unique quadric containing the canonical curve.  Each pencil
has a finite ramification divisor.  Since the marked point varies on the
dominant pointed moduli cover, choose the test fibre so that \(p\) is
unramified for every trigonal pencil, hence for \(J'\). Then

\[h^0(J'(-2p))=0,
\qquad H^0(J(2p))=H^0(J),
\qquad H^0(\mathcal O_C(2p))=H^0(\mathcal O_C).\]

Both equalities of section spaces are induced by the natural inclusions.
If \(z\) is a local parameter at \(p\), then relative to the standard
twisted frames \(z^{-2}e_J\) of \(J(2p)\) and \(z^{-2}\) of
\(\mathcal O_C(2p)\), these inclusions are multiplication by \(z^2\).
After the target-\(J\) bracket block has vanished, every remaining matrix
coefficient of \(\beta\) lies in one of these two images.  Consequently
every coefficient is divisible by \(z^2\), and its reduction on the
contracted fibre is zero:

\[\beta_p=0.\]

Saturation makes \(M_p\ne0\), so \(M_p\) would be a one-dimensional ideal
in the contracted type-\(2+1\) algebra

\[\mathfrak{gl}_2\ltimes
(\mathbb C^2\oplus(\mathbb C^2)^\vee).\]

No such ideal exists: a line in the radical would be a one-dimensional
\(\mathfrak{gl}_2\)-submodule, while a line projecting nontrivially to
\(\mathfrak{gl}_2\) would have to project to its center, whose
nontrivial action on the radical violates ideality. Thus \((2,11)\) is
impossible.

\subsection{\texorpdfstring{The no-high \(q=12\) row is
eliminated}{The no-high q=12 row is eliminated}}\label{the-no-high-q12-row-is-eliminated}

The only row is \((S,q)=(2,12)\), with

\[(\operatorname{rk}G,\deg G,h^0(G))=(7,4,9).\]

The exhaustive equality types are

\[(1,4)+(6,0),\qquad (2,4)+(5,0),\qquad (3,4)+(4,0),\]

and in each case the degree-four factor of rank \(r=1,2,3\) would need
\(r+2\) sections. The degree-four factor lemma rules out all three.
Hence \((2,12)\) is impossible. The \(q=12,13\) zero-weight rows were
already eliminated by the original finite section bounds.

\begin{lemma}[Equality in the BPGN induction]
\label{lem:bpgn-equality-degree}
Let \(C\) be a nonhyperelliptic curve of genus \(g\ge2\), and let \(H\) be
a semistable bundle satisfying

\[0\le\mu(H)\le2g-2,
\qquad h^0(H)=\operatorname{rk}H+\frac{\deg H}{2}.\]

Then \(\deg H\) is divisible by \(2g-2\).
\end{lemma}

\begin{proof}
Follow the induction in the proof of \cite[Theorem~2.1]{BPGN}.  For its
maximal-slope subbundle \(H_1\) and semistable quotient \(H_2\), one has

\[h^0(H)\le h^0(H_1)+h^0(H_2)
\le \operatorname{rk}H+\frac{\deg H}{2}.\]

Equality for \(H\) forces equality throughout, so both \(H_1\) and \(H_2\)
again attain the Clifford bound.  Iteration reduces to line bundles
attaining classical Clifford equality.  No nonspecial term can attain
equality in the stated slope range: Riemann--Roch would force its slope to
be \(2g\).  On a nonhyperelliptic curve, the only line bundles attaining
classical Clifford equality are \(\mathcal O_C\) and \(\omega_C\), of
degrees zero and \(2g-2\).  Additivity of degree along the induction
filtration proves the claim.
\end{proof}

\subsection{\texorpdfstring{The no-high \(q=14\) row is
eliminated}{The no-high q=14 row is eliminated}}\label{the-no-high-q14-row-is-eliminated}

For \((S,q)=(0,14)\) one has

\[(\operatorname{rk}G,\deg G,h^0(G))=(7,8,11).\]

The exact HN enumeration has only one type attaining the BPGN bound
eleven, namely the single factor \((7,8)\); every unstable type has
bound at most ten. Thus \(G\) is semistable (indeed stable, since
\(\gcd(7,8)=1\)) and equality holds in the vector-bundle Clifford
inequality

\[h^0(G)=\operatorname{rk}G+\frac{\deg G}{2}.\]

Apply Lemma~\ref{lem:bpgn-equality-degree}.  In genus four an equality bundle
has degree divisible by \(2g-2=6\), contradicting \(\deg G=8\). Hence
\((0,14)\) is impossible, and \textbf{the complete dense rank-eight
genus-four branch is eliminated}.

\subsection{\texorpdfstring{The high-HN \(q=7\) layer is
eliminated}{The high-HN q=7 layer is eliminated}}\label{the-high-hn-q7-layer-is-eliminated}

The three genus-four high-HN rows with \(q=7\) are

\[(t,y)=(1,5),(2,11),(3,17),\qquad S\ge2.\]

For all three, \(B=A/N\) has rank \(7-t\) and degree \(6(7-t)\). Every
HN slope of \(B\) is at most six, so equality of its mean slope forces
\(B\) to be semistable of slope six. Put \(H=B^\vee\otimes\omega_C\).
Then \(H\) is semistable of degree zero. The exact sequence

\[0\longrightarrow H\longrightarrow G
\longrightarrow N^\vee\otimes\omega_C\longrightarrow0\]

has a negative-slope right term, while \(h^0(G)=4\). Hence \(h^0(H)=4\).
Successively saturating its nowhere-vanishing sections in the semistable
degree-zero bundle gives a saturated evaluation subbundle

\[W=\mathcal O_C^4\subset G.\]

Let \(Q=G/W\). From

\[0\longrightarrow H/W\longrightarrow Q
\longrightarrow N^\vee\otimes\omega_C\longrightarrow0\]

one obtains

\[\mathrm{rk}\,Q=3,\qquad \deg Q=1-S,
\qquad \mu_{\max}(Q)\le0.\]

Parabolic stability gives \(H^0(E)=0\), hence also \(H^0(P)=0\) for the
inverse image \(P\) of \(W\). Since \(\deg M=-1\) and \(h^1(M)=4\), the
inverse-image extension has boundary
\(e:H^0(W)\xrightarrow{\sim}H^1(M)\). Let

\[f:M\otimes W\longrightarrow Q,
\qquad
\phi:H^0(W)\longrightarrow H^0(Q\otimes L)\]

be the quotient bracket and its coefficient map. The same \v{C}ech
comparison gives \(f(v)_*e(u)=f(u)_*e(v)\). If
\(V=\operatorname{im}\phi\) and \(k=\dim V\), the resulting Petri map
lands in \(\mathrm{Sym}^2(V^\vee)\).

Let \(I\subset Q\) be the saturated image of \(f\), with rank \(r\) and
degree \(d\). Every HN slope of \(I\otimes L\) is at most one. The
small-slope section bound gives \(h^0(I\otimes L)\le r\). Since \(V\)
generically spans \(I\otimes L\), one has \(r\le k\le r\), so \(k=r\).

If \(r=2\), then \(-2\le d\le0\). For the quotient line \(R=Q/I\),
Riemann--Roch gives

\[\operatorname{rk}\Theta
=6-d+h^0(Q)-h^0(R)\ge6.\]

Indeed \(h^0(R)\le1\) because \(\deg R\le1\); if it is one, then
\(\deg R\ge0\) forces \(d\le-1\). This contradicts
\(\dim\mathrm{Sym}^2(\mathbb C^2)=3\). If \(r=3\), the Petri map is
injective, but

\[h^0(Q^\vee\otimes\omega_C)=8+S+h^0(Q)\ge10
>6=\dim\mathrm{Sym}^2(\mathbb C^3).\]

Thus \(r=k\le1\). Consequently \(W_0=\ker\phi\) has dimension at least
three. If \(h^0(L)=0\), its remaining self-block vanishes. If
\(L\simeq\mathcal O_C(p)\), write that block as \(sA\). Extension
compatibility gives \(s_*eA=0\), while
\(s_*:H^1(M)\to H^1(\mathcal O_C)\) is an isomorphism by the exact
sequence for \(0\to\mathcal O_C(-p)\to\mathcal O_C\to k(p)\to0\). Hence
\(A=0\). Therefore

\[\mathcal O_C^3\subset\ker\beta.\]

Parabolic stability gives \(H^0(E)=0\), so the full boundary
\(H^0(G)\to H^1(M)\) is injective. The normalizer boundary-cocycle
obstruction excludes the rank-three trivial kernel. This closes all
three \(q=7\) rows. The argument uses closure of the canonical parahoric
lattice under bracket, but no determinant effectivity, puncture-type
exclusion, or relative-section theorem.

\subsection{\texorpdfstring{The high-HN \(q=8\) layer is
eliminated}{The high-HN q=8 layer is eliminated}}\label{the-high-hn-q8-layer-is-eliminated}

The four genus-four high-HN rows with \(q=8\) are

\begin{longtable}[]{@{}rrr@{}}
\toprule\noalign{}
\((t,y)\) & Weight condition & \((\operatorname{rk}H,\deg H)\) \\
\midrule\noalign{}
\endhead
\bottomrule\noalign{}
\endlastfoot
\((1,4)\) & \(S\ge3\) & \((6,0)\) \\
\((1,5)\) & \(S\ge2\) & \((6,1)\) \\
\((2,10)\) & \(S\ge3\) & \((5,0)\) \\
\((2,11)\) & \(S\ge2\) & \((5,1)\) \\
\end{longtable}

Here

\[\deg M=-2,\qquad \deg L=2,\qquad
h^0(G)=h^1(M)=5,\qquad \deg G=2-S,\]

and \(H=(A/N)^\vee\otimes\omega_C\subset G\). The quotient
\(N^\vee\otimes\omega_C\) has negative HN slopes, so \(H^0(H)=H^0(G)\).
We treat the degree-zero and degree-one evaluation types separately.

For \((t,y)=(1,4),(2,10)\), the bundle \(H\) is semistable of degree
zero and its five sections give a saturated

\[W=\mathcal O_C^5\subset G.\]

The quotient \(Q=G/W\) has rank two, degree \(2-S\), and
\(\mu_{\max}(Q)\le0\). Thus \(Q\otimes L\) has determinant of degree at
most three and every line subbundle has degree at most two. If it had
three independent sections, the wedge map to its determinant would have
a decomposable kernel and would produce a degree-at-most-two pencil.
Hence \(h^0(Q\otimes L)\le2\). If a coefficient-rank-two bracket had
saturated image a line \(I\), then \(\mu_{\max}(Q)\le0\) would give
\(\deg(I\otimes L)\le2\), producing the same forbidden pencil. Its image
must therefore have full rank, so the \v{C}ech--Petri map is injective.
But

\[h^0(Q^\vee\otimes\omega_C)=4+S+h^0(Q)\ge7
>3=\dim\operatorname{Sym}^2(\mathbb C^2).\]

The coefficient rank is therefore at most one. The degree-two self-block
can lose at most one further dimension, so three of the five universal
sections lie in the quotient normalizer. The normalizer boundary-cocycle
obstruction eliminates both degree-zero rows.

For \((t,y)=(1,5),(2,11)\), the HN-piece section bounds force the
saturated evaluation bundle to have rank five and degree zero or one.
Degree zero gives \(\mathcal O_C^5\). Degree one gives

\[W=J\oplus\mathcal O_C^4,\qquad J\simeq\mathcal O_C(p),
\qquad F:=\operatorname{im}(\mathcal O_C^5\to G),
\qquad W/F\simeq k(p).\]

When \(\operatorname{rk}H=5\), as in \((2,11)\), only the degree-one
alternative is possible. The saturated \(\mathcal O_C^5\) alternative
occurs only in the row \((1,5)\).

In that saturated alternative, \(Q\) has rank two, degree \(2-S\), and
\(\mu_{\max}(Q)\le1\). If the saturated bracket image is a line and has
two coefficient sections, it has degree one before twisting, so its
twist by \(L\) is a basepoint-free trigonal pencil \(A\). For the
quotient line \(R=Q/I\), the actual restriction image

\[T=\operatorname{im}\bigl(
H^0(Q^\vee\otimes\omega_C)\longrightarrow
H^0(I^\vee\otimes\omega_C)\bigr)\]

has dimension \(2+h^0(Q)\ge2\). \v{C}ech symmetry says that
\(\alpha H^0(A)\) is the same fixed two-plane for every nonzero
\(\alpha\in T\). Two independent choices would have a nonconstant ratio
preserving the pencil; the eigenvector argument makes that ratio
constant, a contradiction.

If the bracket image has full rank, coefficient rank two is excluded by

\[h^0(Q^\vee\otimes\omega_C)=4+S+h^0(Q)\ge6>3.\]

For coefficient ranks greater than two, use

\[0\longrightarrow (H/W)\otimes L\longrightarrow Q\otimes L
\longrightarrow N^\vee\otimes\omega_C\otimes L\longrightarrow0.\]

The two line terms have at most two and one sections, respectively, so
\(h^0(Q\otimes L)\le3\). Thus the only remaining full-image possibility
is coefficient rank three. Its Petri source has dimension at least six.
If the inequality is strict, it cannot inject into
\(\operatorname{Sym}^2(\mathbb C^3)\); in the equality case the map
would be onto, while the evaluation construction supplies a nonzero
reducible conic \(\operatorname{sym}(u_x\otimes\ell_x)\) in its
zero-dimensional orthogonal. Thus this case is impossible as well. The
saturated row again has coefficient rank at most one and yields three
normalizer sections after the one-dimensional self-block loss.

It remains to handle the degree-one saturation defect. Put \(R_0=G/W\).
Then \(R_0\) has rank two, degree \(1-S\), and \(\mu_{\max}(R_0)\le0\).
A rank-one residual bracket image has at most one coefficient section
after twisting by \(L\). A full image has at most two coefficient
sections by the same wedge argument, while coefficient rank two would
inject a space of dimension

\[h^0(R_0^\vee\otimes\omega_C)=5+S+h^0(R_0)\ge7\]

into a three-dimensional symmetric square. Hence the residual
coefficient rank is at most one.

Suppose it is one and write \(\phi(v)=c(v)r\) with \(0\ne r:M\to R_0\).
\v{C}ech symmetry gives \(\dim\ker r_*\ge4\). If \(I\) is the saturation of
\(r(M)\) and \(d=\operatorname{length}(I/r(M))\), then \(0\le d\le2\).
The quotient \(Q_r=R_0/r(M)\) fits into

\[0\longrightarrow I/r(M)\longrightarrow Q_r
\longrightarrow R_0/I\longrightarrow0,
\qquad \deg(R_0/I)=3-S-d.\]

For \(d=0,1,2\), the section bounds for \(R_0/I\) are \(1,1,0\), so
\(h^0(Q_r)\le1,2,2\). The cohomology sequence therefore gives
\(\dim\ker r_*\le2\), a contradiction. Thus the residual coefficient map
vanishes.

Finally, the length-one principal-part map \(W/F\simeq k(p)\) loses at
most one of the five raw sections. If \(L\) is effective, extension
compatibility and the degree-two divisor sequence lose at most one more.
At least three generically independent sections remain in the quotient
normalizer, which is impossible by the boundary-cocycle proposition. The
canonical adjoint parahoric lattice is closed under bracket, so this
calculation remains holomorphic even when the evaluation-defect point is
a puncture.

This eliminates every high-HN \(q=8\) row. The next layer is \(q=9\),
which is also eliminated below; the remaining genus-four high-HN rows
have \(q\ge10\).

\subsection{\texorpdfstring{The high-HN \(q=9\) layer: an apparent
self-block wall and its
closure}{The high-HN q=9 layer: an apparent self-block wall and its closure}}\label{the-high-hn-q9-layer-an-apparent-self-block-wall-and-its-closure}

The four high-HN rows at \(q=9\) are

\begin{longtable}[]{@{}rrr@{}}
\toprule\noalign{}
\((t,y)\) & Weight condition & \((\operatorname{rk}H,\deg H)\) \\
\midrule\noalign{}
\endhead
\bottomrule\noalign{}
\endlastfoot
\((1,3)\) & \(S\ge4\) & \((6,0)\) \\
\((1,4)\) & \(S\ge3\) & \((6,1)\) \\
\((1,5)\) & \(S\ge2\) & \((6,2)\) \\
\((2,11)\) & \(S\ge2\) & \((5,2)\) \\
\end{longtable}

Here again \(H=(A/N)^\vee\otimes\omega_C\), while

\[\deg M=-3,\qquad \deg L=3,\qquad
h^0(G)=h^1(M)=6,\qquad \deg G=3-S.\]

The rank-five row \((2,11)\) is impossible. Every HN factor of \(H\) has
nonnegative degree and the total degree is two. The small-slope bounds
give at most its rank many sections for every degree-zero or degree-one
factor and for every rank-at-least-two degree-two factor; a degree-two
line has at most one section on the nonhyperelliptic fibre. Hence
\(h^0(H)\le5\), contrary to \(h^0(H)=h^0(G)=6\).

For the three \(t=1\) rows, put \(\delta=\deg H=0,1,2\). The six
sections have full-rank raw evaluation

\[F\simeq\mathcal O_C^6\longrightarrow H,
\qquad \operatorname{length}(H/F)=\delta.\]

The residual quotient line \(Q=G/H=N^\vee\otimes\omega_C\) has degree
\(3-S-\delta\), so \(Q\otimes L\) has degree at most two. Its residual
bracket coefficient therefore has rank at most one. If it had rank one,
write it as \(\phi(v)=c(v)r\) with \(0\ne r:M\to Q\). \v{C}ech symmetry and
the boundary isomorphism would give

\[\dim\ker\bigl(H^1(M)\xrightarrow{r_*}H^1(Q)\bigr)\ge5.\]

But \(Q/r(M)\) is torsion of length \(6-S-\delta\le2\), and
\(H^0(Q)=0\). The cohomology sequence gives \(\dim\ker r_*\le2\), a
contradiction. Thus the residual coefficient vanishes in all three rows.

Only the principal parts and the degree-three self-block remain. Killing
the \(\delta\) principal parts leaves at least \(6-\delta\) raw
sections. If \(h^0(L)=0\), the self-block vanishes and the resulting
normalizer ranks are \(6,5,4\); all three rows close. If \(h^0(L)=1\),
multiplication \(H^1(M)\to H^1(\mathcal O_C)\) has a two-dimensional
kernel. The guaranteed normalizer ranks are then

\[(4,3,2)\qquad\text{for }\delta=(0,1,2).\]

The first two rows close by this count. At this stage the row
\((t,y)=(1,5)\) retains a maximal-defect possibility in which both the
principal-part map and the self-block have rank two.

If \(h^0(L)=2\), then \(L\) is a basepoint-free trigonal pencil. The
combined compatibility map from a twelve-dimensional space to
\(H^1(\mathcal O_C)\) is surjective and has kernel dimension eight. As
in the no-high \(q=9\) branch, that kernel admits an injective
six-dimensional matrix pencil, so the extension equations alone force no
constant normalizer kernel. The earliest such wall is

\[(t,y,q)=(1,3,9),\qquad S\ge4,\qquad h^0(L)=2.\]

At the level of the extension equations this is a genuine method wall,
not a geometric adjoint counterexample: generic Kirillov, trace, Jacobi,
and orbit-type identities alone are not enough. The canonical local
lattice supplies the additional input that closes it below.

The sharp primitive pairing-ideal bound above removes every effective
row, including the apparent defect-zero trigonal wall. The quadratic
pushout first gives \(\operatorname{tr}(m^2)=0\), and the symmetric \v{C}ech
identity gives \(P\subset M^\perp\). The global degree comparison then
requires \(\deg D_b\ge S\), whereas the canonical local basis gives
\(\deg D_b\le S/2\). Since every high-HN \(q=9\) row has \(S>0\), this
is a contradiction. Together with the already closed ineffective-\(L\)
subloci and the impossible rank-five row, \textbf{the entire genus-four
high-HN \(q=9\) layer is impossible}. The remaining numerical rows are
treated next.

\subsection{All remaining high-HN rows are
eliminated}\label{all-remaining-high-hn-rows-are-eliminated}

After the preceding \(q\le9\) eliminations, the high-HN ledger contains
only

\begin{longtable}[]{@{}rrr@{}}
\toprule\noalign{}
\(q\) & \((t,y)\) & Weight condition \\
\midrule\noalign{}
\endhead
\bottomrule\noalign{}
\endlastfoot
\(10\) & \((1,4)\) & \(S\ge3\) \\
\(10\) & \((1,5)\) & \(S\ge2\) \\
\(11\) & \((1,5)\) & \(S\ge2\) \\
\end{longtable}

There are no high-HN rows with \(q\ge12\). These last rows have a common
one-line contradiction. Dualizing the contraction sequence and the HN
sequence gives

\[0\longrightarrow M\longrightarrow E\longrightarrow
A^\vee\otimes\omega_C\longrightarrow0,\]

and

\[0\longrightarrow(A/N)^\vee\otimes\omega_C
\longrightarrow A^\vee\otimes\omega_C
\longrightarrow Q:=N^\vee\otimes\omega_C\longrightarrow0.\]

Thus \(Q\) is a genuine quotient line of \(E\), with no torsion
correction, and

\[\deg Q=6-S-y.\]

Let \(a\) and \(b\) denote the numbers of type-\(2+1\) and full-flag
nonscalar markings. Then \(S=2a+3b\). Scalar markings give quotient
weight zero, while the induced quotient-line weight at every nonscalar
marking lies in \([0,1)\). If \(w_Q\) is the total quotient weight, then

\[w_Q<a+b,
\qquad
\operatorname{pardeg}Q
<6-y-S+a+b=6-y-a-2b.\]

For \(y=4\) the condition \(S\ge3\) implies \(a+2b\ge2\), so
\(\operatorname{pardeg}Q<0\). For \(y=5\) the condition \(S\ge2\)
implies \(a+2b\ge1\), and again \(\operatorname{pardeg}Q<0\). Parabolic
stability of the degree-zero canonical adjoint bundle requires every
proper quotient line to have strictly positive parabolic degree. Both
possibilities are therefore impossible. Consequently \textbf{every
genus-four high-HN row is eliminated}.

\section{Proper projective-closure coefficients in genus
four}\label{proper-projective-closure-coefficients-in-genus-four}

The dense rank-eight calculation is not by itself enough. The projective
closure alternatives in the companion manuscript \cite{SneidermanG5} also
contain

\begin{longtable}[]{@{}ll@{}}
\toprule\noalign{}
Projective closure & Adjoint constituents \\
\midrule\noalign{}
\endhead
\bottomrule\noalign{}
\endlastfoot
\(\mathrm{PGL}_2\simeq\mathrm{SO}_3\) & \(U_5\oplus U_3\) \\
monomial with component quotient \(S_3\) & \(U_6\oplus U_2\) \\
monomial with component quotient \(C_3\) &
\(U_3^+\oplus U_3^-\oplus\chi\oplus\chi^{-1}\) \\
\end{longtable}

The rank-three, rank-two, and rank-one constituents are already below
the fixed-part threshold in genus four. It remains to exclude a rank-one
fixed part for \(U_5\) and \(U_6\).

\subsection{The symmetric universal-extension
calculation}\label{the-symmetric-universal-extension-calculation}

Both remaining coefficients carry an invariant symmetric form

\[B:E\otimes E\longrightarrow\mathcal O_C\]

which is holomorphic on the canonical Deligne lattice and nondegenerate
over the function field. Reciprocal nonzero parabolic weights pair with
one zero of \(B\) at a marking; this creates no pole or hidden
\(\mathcal O_C(D)\) target.

Suppose an inverse image \(P\subseteq E\) fits into

\[0\longrightarrow M\longrightarrow P
\longrightarrow V\otimes\mathcal O_C\longrightarrow0\]

and its boundary \(e:V\to H^1(M)\) is an isomorphism. Put \(L=M^{-1}\).
The restriction of the pairing is a morphism

\[Q_2=B|_{M\otimes M}:M\longrightarrow L.\]

Pairing \(P\) against \(M\) extends \(Q_2\), so the pushout of the
extension along \(Q_2\) splits and

\[(Q_2)_*e=0.\]

Whenever \(H^0(\omega_C\otimes L^{-1})\ne0\), a nonzero \(Q_2\) would
have injective Serre-dual multiplication

\[H^0(\omega_C\otimes L^{-1})
\xrightarrow{\,Q_2\,}H^0(\omega_C\otimes L),\]

contradicting \((Q_2)_*=0\) and surjectivity of \(e\). Thus \(Q_2=0\) in
all cases below. Pairing with \(M\) then descends on the trivial
quotient to a map \(c:V\otimes\mathcal O_C\to L\). Comparing the pairing
in local splittings of the extension gives the symmetric \v{C}ech identity

\[c(v)_*e(u)+c(u)_*e(v)=0
\quad\text{in }H^1(\mathcal O_C).\]

This identity will force \(c=0\), or will identify the inverse image
\(P\) with \(M^\perp\).

\subsection{The rank-five orthogonal
coefficient}\label{the-rank-five-orthogonal-coefficient}

Here \(d=5=g+1\). The companion manuscript's first-super
proposition \cite{SneidermanG5} excludes every total weight \(S\ne1\). Its weight-one
classification leaves exactly

\[0\longrightarrow M\longrightarrow E_5
\longrightarrow\mathcal O_C^4\longrightarrow0,
\qquad \deg M=-1,\]

with \(e:\mathbb C^4\simeq H^1(M)\). The local \(U_5\) weights are
\(0,\pm\alpha,\pm2\alpha\); total one occurs at a projective involution,
with canonical weights \((0,0,0,\tfrac12,\tfrac12)\).

The symmetric calculation gives \(Q_2=0\): Riemann--Roch gives

\[h^0(\omega_C\otimes L^{-1})=2+h^0(L)\ge2.\]

Since \(\deg L=1\), one has \(h^0(L)\le1\). If it is zero, the descended
map \(c:\mathcal O_C^4\to L\) vanishes. If \(L\simeq\mathcal O_C(p)\)
with section \(s\), then

\[s_*:H^1(M)\longrightarrow H^1(\mathcal O_C)\]

is an isomorphism by \(0\to M\xrightarrow{s}\mathcal O_C\to k(p)\to0\).
The symmetric \v{C}ech identity therefore forces \(c=0\) again. Thus \(M\)
lies in the generic radical of \(B\), contradicting generic
nondegeneracy. Hence the rank-five constituent has no rank-one fixed
part.

\subsection{Exact rank-six monomial
arithmetic}\label{exact-rank-six-monomial-arithmetic}

For \(U_6\), the positive local total weights are

\[\frac32,\quad2,\quad\frac52,\quad3,\]

and the global total \(S\) is integral. If a rank-zero Hodge vector
existed, put

\[F=E_6^\vee\otimes\omega_C,
\qquad \deg F=36+S,
\qquad h^0(F)=18+S,\]

and let \(A\subset F\) be its rank-five contraction kernel. Write
\(q=\deg(F/A)=6-\deg M\ge7\).

In the no-high case the exact inequalities are

\[S\le\frac{24}{7},
\qquad
\max(7,6+S)\le q\le10-S.\]

HN-piece section counting gives

\begin{longtable}[]{@{}rrrrr@{}}
\toprule\noalign{}
\(S\) & \(q\) & \(\deg(E_6/M)\) & \(h^0(E_6/M)\) & HN upper bound \\
\midrule\noalign{}
\endhead
\bottomrule\noalign{}
\endlastfoot
\(0\) & \(7\) & \(1\) & \(4\) & \(4\) \\
\(0\) & \(8\) & \(2\) & \(5\) & \(4\) \\
\(0\) & \(9\) & \(3\) & \(6\) & \(4\) \\
\(0\) & \(10\) & \(4\) & \(7\) & \(5\) \\
\(2\) & \(8\) & \(0\) & \(5\) & \(5\) \\
\end{longtable}

Thus only \((S,q)=(0,7)\) and \((2,8)\) remain. In the high-HN case the
exact offsets are

\[t=1:\ y=5,6,
\qquad
t=2:\ y=12,\]

where \(y=\deg N-S\). At the top offsets \(y=6,12\), the genuine
quotient \(Q=N^\vee\otimes\omega_C\) has ordinary degree \(-S\). Its
induced quotient weight is at most the total coefficient weight \(S\),
so \(\operatorname{pardeg}Q\le0\), contradicting stability. The sole
strict row is

\[(t,y,q)=(1,5,7),\qquad S\ge2.\]

Here \(Q=N^\vee\otimes\omega_C\) is a genuine quotient line of \(E_6\)
with \(\deg Q=1-S\). If \(m\) is the number of nonscalar markings, the
local total list gives

\[m\le\left\lfloor\frac{2S}{3}\right\rfloor\le S-1.\]

The quotient weight is strictly below \(m\), so

\[\operatorname{pardeg}Q<1-S+m\le0,\]

contradicting parabolic stability. No high-HN row survives.

\subsection{The two rank-six no-high
rows}\label{the-two-rank-six-no-high-rows}

For \((S,q)=(2,8)\), equality gives

\[E_6/M\simeq\mathcal O_C^5,
\qquad \deg M=-2,
\qquad \mathbb C^5\simeq H^1(M).\]

The symmetric calculation gives \(Q_2=0\), since
\(h^0(\omega_C\otimes L^{-1})=1+h^0(L)>0\). On the nonhyperelliptic test
fibre \(h^0(L)\le1\). If \(L\) is effective, multiplication
\(H^1(M)\to H^1(\mathcal O_C)\) has one-dimensional kernel, whereas a
nonzero cross term would put a four-dimensional subspace into that
kernel. Thus the cross term vanishes. Since the quotient is globally
generated, \(M\) would be in the generic radical of \(B\), a
contradiction.

For \((S,q)=(0,7)\), let \(R\subset G=E_6/M\) be proper. Its inverse
image in \(E_6\) has degree \(\deg R-1\), so stability of the ordinary
degree-zero coefficient gives \(\deg R\le0\). Hence \(G\) is stable of
rank five and degree one. The saturation of its four-section evaluation
is generically generated and has degree zero; all four sections force
its rank to be four. It is therefore

\[W=\mathcal O_C^4\subset G.\]

Let \(P\subset E_6\) be its inverse image. Then \(H^0(P)=0\) and the
boundary \(\mathbb C^4\simeq H^1(M)\) is an isomorphism. The symmetric
calculation gives \(Q_2=0\) and kills the cross term on \(W\), so
\(P\subseteq M^\perp\). Both bundles have rank five and degree minus
one; hence

\[P=M^\perp,
\qquad M^\perp/M\simeq\mathcal O_C^4,
\qquad \det E_6\simeq\mathcal O_C.\]

This contradicts the determinant of the actual \(S_3\) monomial
constituent. The diagonal torus factors on the six ordered roots
multiply to one, while a transposition permutes them by three swaps.
Consequently

\[\det U_6=\operatorname{sgn}:S_3\longrightarrow\{\pm1\}.\]

The component quotient is \(S_3\), so density makes this a nontrivial
unitary determinant character. Its unpunctured holomorphic line cannot
be trivial. This closes the last proper-closure coefficient.

\section{Genus-four propagation}\label{genus-four-propagation}

The preceding fixed-part vanishing holds on every connected finite-etale
cover used in the companion manuscript \cite{SneidermanG5}. Projective
finite-index descent and direct
Artin devissage have no further genus restriction. The reducible
nonscalar multiplicity comparisons remain strict at \(g=4\):

\[(u,w)=(1,\le4),(2,\le3),(3,1),
\qquad
w<2g-2u=6,4,2.\]

\begin{lemma}[The standard subsystem on a finite-index base]
\label{lem:finite-index-standard}
Let \(B\to\mathcal M_{g,n}\) be a connected base whose mapping-class-group
image has finite index, with labeled marked sections, and put

\[H=R^1\pi_*\mathbb C\subset R^1\pi_*^\circ\mathbb C.\]

Then \(H\) is irreducible with Zariski-dense \(\Sp_{2g}\)-monodromy, and
the quotient \(R^1\pi_*^\circ\mathbb C/H\) is constant.  Consequently
every irreducible rank-\(2g\) sub-local system of
\(R^1\pi_*^\circ\mathbb C\) equals \(H\), and \(H\) has no nonzero
invariant symmetric bilinear form.
\end{lemma}

\begin{proof}
The residue sequence identifies the quotient with the constant local system

\[\ker\!\left(\mathbb C^n\xrightarrow{\sum}\mathbb C\right)(-1).\]

The standard action of \(\Mod_{g,n}\) on \(H\) has image
\(\Sp_{2g}(\mathbb Z)\).  The image of a finite-index subgroup is again
finite index, hence Zariski dense in \(\Sp_{2g}\); the standard
representation is therefore irreducible.

Let \(I\) be an irreducible rank-\(2g\) sub-local system of the punctured
cohomology.  If \(I\cap H=0\), then \(I\) injects into the constant residue
quotient.  Its monodromy is then trivial, contradicting irreducibility and
rank \(2g>1\).  Thus \(I\cap H\ne0\); irreducibility and equality of ranks
give \(I=H\).  Finally, the standard \(\Sp_{2g}\)-representation preserves
its alternating form and has no nonzero invariant symmetric form.
\end{proof}

The scalar residual case needs one new equality argument. The companion
manuscript's exact trivial-coefficient gap \cite{SneidermanG5} says every nonzero
sub-local system of \(R^1\pi_*^\circ\mathbb C\) has rank at least
\(2g=8\). A first nonzero scalar Artin invariant has image \(I\) of rank
at most eight, hence exactly eight and irreducible. By
Lemma~\ref{lem:finite-index-standard},

\[I=R^1\pi_*\mathbb C,
\qquad (W^0)^\vee\simeq R^1\pi_*\mathbb C.\]

The left side is an adjoint \(\mathfrak{sl}_3\) local system and therefore
preserves its nondegenerate symmetric Killing form. The right side has
Zariski-dense \(\Sp_8\)-monodromy and admits no invariant symmetric form.
This contradiction closes the scalar equality case.

The remaining propagation margins are still strict. The isomonodromy
range is \(3<2\sqrt5\), and for a characteristic extension of total rank
at most three,

\[n_i d_i\le2<2g-2d_i,
\qquad d_i\le2,\]

because the right side is at least four. The projective/linear
integrality, unitary-conjugate, and socle-induction arguments therefore
carry over with the same cited inputs.

\begin{proof}[Completion of Candidate theorem~\ref{cand:main}]
A reducible semisimple rank-three representation has constituents of rank at
most two, hence lies in the published strict range.  For an irreducible
system with infinite projective image, the imported projective-closure
classification leaves the four rows in \cref{sec:interfaces}.  All
constituents of rank at most three are already below the fixed-part
threshold.  The dense coefficient $U_8$ is eliminated by the no-high and
high-HN arguments above; the proper-closure coefficients $U_5$ and $U_6$
are eliminated in
\cref{proper-projective-closure-coefficients-in-genus-four}.  Thus the
required finite-cover adjoint fixed-part vanishing holds for every actual
irreducible coefficient.

Projective finite-index descent and direct Artin devissage import that
vanishing to the dominant versal base.  Reducible nonscalar residual systems
retain strict multiplicity margins.  The only scalar equality is excluded in
\cref{genus-four-propagation} by the incompatibility between an adjoint
symmetric Killing form and the standard $\Sp_8$ representation.  The
remaining rigidity, integrality, unitary-conjugate, and semisimple
nonabelian-Hodge steps are B5.  Finally, for a characteristic extension of
total rank at most three, $n_i d_i\le2<2g-2d_i$ for $d_i\le2$ at $g=4$;
the imported socle induction therefore handles the nonsemisimple case.
This proves the stated conclusion conditional on B1--B5.
\end{proof}

\section{Status and remaining verification boundary}
\label{sec:status}

The status of Candidate theorem~\ref{cand:main} is \texttt{CANDIDATE}.  More precisely:
\begin{itemize}
\item The finite HN ledgers and the displayed dimension, determinant, and
fibrewise Lie-algebra calculations are exact finite calculations recorded in
this manuscript.
\item The dense $U_8$ elimination is a proposed proof conditional on the
fixed-part, canonical-lattice, parabolic, and extension interfaces B1--B4.
The most sensitive new steps are the \v{C}ech--Petri symmetry, saturation of
raw evaluation lattices, the Spin lift and parity comparison, the primitive
pairing-ideal bound at punctures, and the equality analysis in the BPGN
Clifford induction.
\item The $U_5$ and $U_6$ eliminations additionally depend on the companion
first-super classification and on the invariant symmetric form extending
holomorphically on the canonical Deligne lattice.
\item Formal propagation at scalar residual type uses the asserted
Zariski-dense $\Sp_8$ monodromy on every finite-index mapping-class base and
the absence of an invariant symmetric form on its standard representation.
\item The conclusion also relies on B5: rigidity, integrality,
isomonodromy, finite-cover descent, and socle induction.  Those steps are not
reproved here.
\end{itemize}

Independent specialist reconstruction of these geometric steps remains
necessary.  Finite scripts check only enumerations and selected linear
algebra; the TeX build checks only the document.  The manuscript includes
only arguments used in the proposed proof.
No external referee acceptance, priority claim, or established theorem
status is asserted.

\paragraph{Assistance disclosure.}
OpenAI Codex (GPT-5.6) assisted with preparation and internal checking of this
manuscript.  The author reviewed the manuscript and is responsible for all
statements, proofs, and errors.

{\small
\begin{thebibliography}{99}

\bibitem{BPGN}
L.~Brambila-Paz, I.~Grzegorczyk, and P.~E. Newstead,
\emph{Geography of Brill--Noether loci for small slopes},
J. Algebraic Geom. \textbf{6} (1997), 645--669.

\bibitem{ChenSalter}
L.~Chen and N.~Salter,
\emph{The Birman exact sequence does not virtually split},
Math. Res. Lett. \textbf{28} (2021), 383--413.

\bibitem{EarleKra}
C.~J. Earle and I.~Kra,
\emph{On sections of some holomorphic families of closed Riemann surfaces},
Acta Math. \textbf{137} (1976), 49--79.

\bibitem{LLCan}
A.~Landesman and D.~Litt,
\emph{Canonical representations of surface groups},
Ann. of Math. (2) \textbf{199} (2024), 823--897.

\bibitem{LLGeo}
A.~Landesman and D.~Litt,
\emph{Geometric local systems on very general curves and isomonodromy},
J. Amer. Math. Soc. \textbf{37} (2024), 683--729.

\bibitem{SneidermanG5}
R.~Sneiderman,
\emph{Rank-Three Finite-Image Theorem for Surface Groups in Genus at Least
Five}, companion candidate manuscript, version 0.1.5-candidate, 19 August
2026, \url{https://github.com/Robby955/rank-three-mcg-finiteness}.

\end{thebibliography}
}

\end{document}
